\documentclass[12pt, a4paper, oneside]{report}

% ,  PACKAGES , 
\usepackage[margin=2cm]{geometry}
\usepackage[utf8]{inputenc}
\usepackage{amsmath, amssymb, amsthm, bm}
\usepackage{graphicx}
\usepackage[colorlinks=true, linkcolor=blue, citecolor=blue, urlcolor=blue]{hyperref}
\usepackage{booktabs}
\usepackage{setspace}
\usepackage{caption}
\usepackage{subcaption}
\usepackage{algorithm}
\usepackage{algpseudocode}
\usepackage{physics}
\usepackage{qcircuit}
\usepackage{tcolorbox}
\usepackage{float}
\usepackage{cite}
\onehalfspacing

\begin{document}

% ==========================================
% TITLE PAGE
% ==========================================
\begin{titlepage}
    \centering
    \vspace*{1cm}
    \includegraphics[width=0.5\textwidth]{imperial_logo.png} \\[1.5cm]
    {\large \textbf{Project Code: SPC-Appelbe-1}} \\[0.5cm]
    {\huge \bfseries Quantum-Assisted Adaptive Mesh Refinement for MHD} \\[0.5cm]
    {\Large \bfseries A Hybrid QAOA Framework and its Structural Limits for Topological Instability Detection} \\[2cm]

    \textbf{Student CID:} 06068471 \\[0.5cm]
    \textbf{Supervisor:} Dr.\ Brian Appelbe \\
    \textbf{Assessor:} Dr.\ Aidan Crilly \\[2cm]

    \textbf{Word Count}: 10 085 \\

    \vfill
    Academic Year 2025--2026
\end{titlepage}

% ==========================================
% LAYPERSON'S SUMMARY (commented out for arXiv version; keep for thesis)
% ==========================================
\iffalse
\chapter*{Layperson’s Summary}
\addcontentsline{toc}{chapter}{Layperson’s Summary}

\begin{center}
{\Large \textbf{Steering the Storm: Using Quantum Ideas to Anticipate Plasma Instabilities}}
\end{center}

\noindent Fusion energy, the same process that powers the Sun, promises virtually limitless, clean power. To achieve it on Earth, scientists must confine plasma, a gas so hot that its atoms break apart into electrically charged particles, inside powerful magnetic fields. The challenge is that plasma is inherently unruly: it twists, flows, and can suddenly become unstable. Small disturbances can cascade into larger disruptions capable of damaging the walls of a reactor. Predicting when and where these instabilities will emerge is one of the central challenges in plasma physics.

To study plasma, scientists run large-scale computer simulations. The plasma volume is divided into millions of small cells, and the equations governing its motion are solved in every cell at each instant. Smaller cells yield more accurate results, but also demand far more computing power. Using the finest resolution everywhere would require months of supercomputer time for a single simulation.

A common solution is \textit{Adaptive Mesh Refinement} (AMR): the computer increases resolution only where complex behaviour is detected. However, current AMR methods are too reactive, they refine the grid only after an instability has already grown large enough to be noticed. By that point, it is possible that the simulation have already lost important detail.

This project asks a different question: rather than reacting to instability, can we anticipate it? The approach investigated here uses a small quantum circuit as an early-warning system. Quantum computers process information using quantum bits (qubits), which can represent multiple possibilities at once through a property called \textit{superposition}. Another property, \textit{interference}, allows qubits to amplify useful signals while cancelling out noise (like the way noise-cancelling headphones work). These two properties, used in the right way, can provide huge advantage over classical methods.

In the Q-HAS (Quantum-Hierarchical Adaptive Steering) framework developed in this project, indicators of plasma behaviour, such as the rate of change of flow velocity and magnetic field, are encoded into a small set of qubits. A quantum algorithm then evaluates whether the pattern of encoded signals is consistent with the early stages of a known instability, such as a shearing flow, a forming vortex, or a magnetic reconnection event. The quantum circuit does not directly solve the plasma equations; instead, it only guides the classical simulation by flagging regions that are likely to become unstable in the near future.

The result is a hybrid architecture: a classical supercomputer performs the detailed plasma calculations, while the quantum component acts as a predictive steering mechanism that focuses computational effort where it is most needed. Both a theoretical analysis and a proof-of-concept software implementation are presented, demonstrating the framework on eight standard plasma test cases including the Kelvin-Helmholtz instability, the Orszag-Tang vortex, Harris current sheets, and isolated vortex, reconnection, and rotor scenarios.

Although today’s quantum hardware remains limited in scale, this work provides a foundation for integrating quantum sensing into classical simulations. It illustrates that quantum computing may be most valuable not as a replacement for classical methods, but as a specialised complement that detects an instability before it becomes critical.
\fi

% ==========================================
% ABSTRACT 
% ==========================================
\begin{abstract}
\noindent \textbf{Background.} Adaptive Mesh Refinement (AMR) is standard for multi-scale magnetohydrodynamic (MHD) simulations, but conventional gradient-based refinement criteria are reactive: they trigger only once an instability has reached detectable amplitude. Whether quantum variational subroutines can provide a predictive, topology-aware refinement criterion is an open question.

\noindent \textbf{Contribution.} We introduce Q-HAS, a hybrid quantum-classical framework that reformulates instability detection as the minimisation of a structured local Ising Hamiltonian on a coarse-grained flux graph. Each qubit's polar angle $\theta$ encodes a classical multi-indicator score, while its azimuthal phase $\phi$ encodes the temporal growth rate of stress flux. A QAOA-based algorithm evaluates activity-bias ($Z$), uncertainty-weighted gradient-coupling ($ZZ$) and circulation-plaquette ($ZZZZ$) terms under a decoupled $f \times g \times T_{\mathrm{rcf}}$ coefficient architecture. We derive a null-gradient result showing that optimisation requires mixer-induced symmetry breaking, together with explicit gradient expressions linking plasma observables to qubit rotations.

\noindent \textbf{Empirical results ($N = 256$, $\mathrm{Re} = \mathrm{Rm} = 400$, 8 qubits, Qiskit-Aer).} Per-scenario training yields a $0.66\%$ composite-loss advantage over a multi-indicator classical baseline, with within-scenario gains on Kelvin--Helmholtz and Harris tearing and no advantage (or regressions) on Orszag--Tang and MHD rotor.

\noindent \textbf{Falsification study.} A leave-one-scenario-out (LOSO) evaluation, together with machine-learning surrogates bounding the entire family of per-site decision functions over the 9 physical features used, reveals a structural limit: the mean-field ceiling collapses from $F_1 = 0.986$ (random split) to $F_1 = 0.189$ (LOSO), well below the classical baseline ($F_1 = 0.434$). A formal proposition establishes that, within this feature basis and for couplings up to nearest-neighbour stencil range, no exact ground state of a local Ising Hamiltonian can exceed this empirical bound by more than $\lesssim 0.01$, regardless of solver, circuit depth, or noise. The bottleneck is feature insufficiency on held-out instability classes, not the quantum optimiser.

\noindent \textbf{Scope.} Conclusions apply to 2D incompressible MHD in the laminar-to-transitional regime ($\mathrm{Re} \le 1600$) tested here. All circuits are executed on the Qiskit-Aer statevector simulator; hardware execution is not reported.
\end{abstract}

\tableofcontents

\chapter{Introduction}
\label{chap:introduction}

Magnetohydrodynamics (MHD)~\cite{freidberg2014,goedbloed2004} governs the nonlinear coupling of fluid motion, electromagnetic fields, pressure gradients, and topological constraints in electrically conducting plasmas. In regimes relevant to magnetic confinement fusion~\cite{wesson2011}, space plasmas, and astrophysical flows, small perturbations evolve into coherent instabilities: vortices, current sheets, magnetic reconnection events, and kink modes. Predicting and controlling these structures is a central challenge in computational plasma physics.

The MHD equations combine the Navier--Stokes equations with Maxwell's equations. In incompressible MHD (relevant to tokamak interiors), the magnetic field evolves through advection and resistive diffusion. The nonlinear coupling between velocity and magnetic fields produces instabilities: shearing flows trigger the Kelvin-Helmholtz instability~\cite{chandrasekhar1961}, current sheets undergo tearing-mode reconnection~\cite{furth1963}, and magnetic flux tubes develop helical kink modes~\cite{bateman1978}.

Analytical solutions are unavailable in realistic regimes; progress depends on numerical simulation. Plasma dynamics are multiscale: fine-scale structures coexist with large-scale flows, and uniform resolution is prohibitive. Direct Numerical Simulation (DNS) cost scales as $Re^{9/4}$ in three dimensions~\cite{pope2000}, becoming intractable at high Reynolds numbers $Re$ and magnetic Reynolds numbers $Rm$. Adaptive Mesh Refinement (AMR)~\cite{berger1984,berger1989} increases resolution only where strong gradients or discontinuities are detected. However, current AMR refinement criteria are local and reactive: they respond only after instabilities have reached detectable amplitude.

Quantum computing offers alternative paradigms based on superposition, interference, and entanglement~\cite{nielsen2010}. Quantum approaches to linear algebra (HHL algorithm~\cite{harrow2009}) and quantum many-body simulation~\cite{feynman1982} are well studied, but direct quantum solution of nonlinear PDEs remains challenging: quantum circuits evolve linearly and unitarily, whereas MHD systems are nonlinear. Rather than forcing the two mathematical frameworks to reconcile, we shift our focus toward AMR as a strategic alternative to direct integration. Hybrid strategies combining classical solvers with targeted quantum subroutines offer a more realistic near-term pathway.

The Quantum Approximate Optimisation Algorithm (QAOA)~\cite{farhi2014} alternates between a ``cost'' unitary (encoding problem structure) and a ``mixer'' unitary (generating quantum fluctuations), with variational parameters optimised by a classical loop. Its fixed-depth circuit structure and noise resilience make it suitable for Noisy Intermediate-Scale Quantum (NISQ) devices~\cite{preskill2018}.

This project develops a quantum-assisted steering mechanism for AMR: the Quantum-Hierarchical Adaptive Steering framework (Q-HAS). The quantum component serves a diagnostic role rather than replacing the classical solver. Instability detection is reformulated as the minimisation of a structured (quantum Ising) Hamiltonian defined on a coarse-grained flux representation of the plasma domain. Physical quantities (flux amplitudes and temporal variations) are encoded into qubit states via Bloch sphere mapping, and a QAOA-based variational algorithm evaluates a cost function whose terms correspond to instability signatures like activity bias, spatial gradient coupling, physical circulation, and X-point reconnection.

The objectives of the project are threefold:
\begin{enumerate}
    \item To construct a consistent mapping between coarse-grained MHD flux variables and qubit representations, encoding both the amplitude and the growth rate of local plasma quantities.
    \item To define and analyse the Hamiltonian governing instability detection, including the behaviour of its variational gradients and the role of mixer-induced symmetry breaking.
    \item To implement the framework in software and assess, both theoretically and computationally, how such a hybrid mechanism can reduce unnecessary refinement and improve simulation efficiency.
\end{enumerate}

\section*{Related Work}
\addcontentsline{toc}{section}{Related Work}

\paragraph{Quantum algorithms for fluid and plasma simulation.}
Direct quantum approaches to nonlinear PDEs include quantum algorithms for the Navier--Stokes equations~\cite{gaitan2020}, Carleman linearisation techniques for fluid-like systems~\cite{liu2021carleman}, and general results on the hardness of simulating nonlinear dynamics unitarily~\cite{childs2020pde,dodin2021}. Quantum signal processing and block-encoded linear solvers (QSVT, quantum linear systems) are being explored as PDE subroutines~\cite{martyn2021qsvt}. In contrast to these ``replace the solver'' approaches, the present work follows the complementary direction of using a quantum subroutine as a \emph{diagnostic} rather than an integrator, which is closer in spirit to recent proposals for quantum-assisted refinement heuristics and to classical machine-learning approaches to AMR and disruption prediction in tokamaks~\cite{katesharbeck2019}.

\paragraph{Variational quantum algorithms and warm-starting.}
The QAOA~\cite{farhi2014} and its Quantum Alternating Operator Ansatz generalisation~\cite{hadfield2019} form the theoretical backbone of the present framework. We draw in particular on warm-start QAOA formulations~\cite{egger2021warmstart}, which initialise the variational state close to a classical solution and help mitigate the barren plateau phenomenon~\cite{mcclean2018,cerezo2021cost}. Studies of parameter concentration and transferability in QAOA~\cite{zhou2020qaoa,wang2018qaoa} motivate our use of a fixed linear-ramp schedule for the structural angles.

\paragraph{Expressivity and learnability of variational models.}
The structural question we raise in Chapter~\ref{chap:falsification} --- whether a family of local Ising Hamiltonians can represent a given decision function --- is the applied counterpart of a growing literature on the expressivity and learnability of variational quantum models~\cite{schuld2021effect,cerezo2022vqa,du2022learnability}. Our ``mean-field ceiling'' proposition is close in spirit to no-go results showing that shallow local circuits cannot exceed certain classical surrogates on specific tasks.

\paragraph{Adaptive Mesh Refinement and MHD codes.}
The AMR component of this work builds on the classical Berger--Oliger--Colella framework~\cite{berger1984,berger1989} and the indicator-based refinement criteria used in production MHD codes such as FLASH, Athena and PLUTO~\cite{fryxell2000,stone2008,mignone2007}. The L\"ohner second-derivative error estimator~\cite{lohner1987} and the Okubo--Weiss $Q$ criterion~\cite{okubo1970,weiss1991} enter directly in our classical multi-indicator score.

\section*{Contributions and scope}
\addcontentsline{toc}{section}{Contributions and scope}

The specific contributions of this work are:
\begin{enumerate}
    \item A physically motivated encoding of coarse-grained MHD observables into a local Ising Hamiltonian, combining amplitude ($\theta$) and temporal growth ($\phi$) in a Bloch-sphere representation, with a decoupled $f \times g \times T_{\mathrm{rcf}}$ coefficient architecture (Chapters~\ref{chap:flux-graph}--\ref{chap:hamiltonian}).
    \item A null-gradient result identifying mixer-induced symmetry breaking as the ignition mechanism of the variational optimisation, together with explicit gradient expressions connecting plasma observables to Bloch-sphere torques (Chapter~\ref{chap:theory}).
    \item A full software implementation with a 4th-order finite-difference MHD solver, a Qiskit VQA pipeline, Optuna-based Bayesian training and a multi-tier validation suite (Chapter~\ref{chap:implementation}), together with empirical per-scenario benchmarks (Chapter~\ref{chap:results}).
    \item A cross-scenario falsification study, including a formal ceiling proposition bounding the $F_1$ of any local Ising Hamiltonian over the 9-feature per-site basis used here, irrespective of solver, depth, or noise (Chapter~\ref{chap:falsification}).
\end{enumerate}

\paragraph{Scope and caveats.}
All results are obtained in 2D incompressible MHD at $\mathrm{Re} = \mathrm{Rm} = 400$ (with scaling checks up to $\mathrm{Re} = 1600$), grid size $N = 256$, and 8-qubit QAOA circuits executed on the Qiskit-Aer statevector simulator. Neither quantum hardware execution nor noisy simulation is used for the headline results. The classical baseline is a multi-indicator threshold AMR implemented in the same code base, not an external production code. These restrictions are revisited in the Discussion (Chapter~\ref{chap:discussion}).

The complete source code is publicly available at \url{https://github.com/armandld/BA_Proj}.

\chapter{Problem Formulation: The Flux Graph}
\label{chap:flux-graph}

Q-HAS represents the plasma domain as a graph whose edges carry physical flux information. This chapter defines the coarse-grained flux representation and the mapping from classical plasma quantities to qubit states.

\section{The Dual Graph Representation}

Let $\Omega \subset \mathbb{R}^3$ be the plasma volume. We decompose $\Omega$ into a set of coarse blocks $\mathcal{B} = \{b_1, \dots, b_N\}$.
Instead of nodal values, we define the state of the system by the flux exchange between these blocks.

Let $G = (V, E)$ be the dual graph where:
\begin{itemize}
    \item $V$ represents the volume blocks (nodes).
    \item $E$ represents the interfaces (surfaces or fluxes) between adjacent blocks.
\end{itemize}

We map the physical flux $\Phi_{ij}$ through each edge $(i,j) \in E$ to a qubit state.

\section{Classical Flux Extraction}

Let $U_i(t)$ and $U_j(t)$ be the classical state vectors at nodes $i$ and $j$ at time $t$. The implementation extracts five normalised MHD fields: $\tilde{v}_x = v_x / v_0$, $\tilde{v}_y = v_y / v_0$, $\tilde{B}_x = B_x / B_0$, $\tilde{B}_y = B_y / B_0$, and $\tilde{J}_z = J_z / J_0$, where $v_0$, $B_0$, and $J_0$ are characteristic reference scales. When the considered block is coarse, the maximum value of these fields within the block is used.

For each edge (in 2D here), the differences between neighbouring cells are computed in both directions (right/left for horizontal edges, top/bottom for vertical edges). These differences are then processed through a shock-diode logic that separates the physically distinct contributions of compression and shear:

\begin{itemize}
    \item \textbf{Shock component (compression only):} Only negative normal differences (compression) are retained: $\Phi_{\text{shock}} = \sqrt{\max(0, -\Delta v_n)^2 + \max(0, -\Delta B_n)^2}$. Expansion (positive $\Delta v_n$) is filtered out as physically "harmless".
    \item \textbf{Shear component (always dangerous):} Tangential differences are taken in absolute value regardless of sign: $\Phi_{\text{shear}} = \sqrt{|\Delta v_t|^2 + |\Delta B_t|^2}$. Sliding motion is unstable in either direction.
    \item \textbf{Weighted combination:} The total stress flux combines both contributions with tuneable weights:
    \begin{equation}
        \Phi_{ij,l}(t) = \sqrt{(w_{\text{shock}} \cdot \Phi_{\text{shock}})^2 + (w_{\text{shear}} \cdot \Phi_{\text{shear}})^2 + (w_J \cdot \Delta J_z)^2}
    \end{equation}
\end{itemize}

The maximum danger from either side of the edge is retained: $\Phi_{ij} = \max(\Phi_{ij}^{\text{right}}, \Phi_{ij}^{\text{left}})$.

The time derivative $\dot{\Phi}_{ij}(t)$, representing the flux evolution rate, is computed as the difference between consecutive VQA updates:
\begin{equation}
    \dot{\Phi}_{ij, l}(t) = \Phi_{ij, l}(t) - \Phi_{ij, l}(t - \Delta t_{\text{hybrid}})
\end{equation}
Note that no explicit $\Delta t$ division is needed: both the numerator and denominator of the phase encoding (Equation~\ref{eq:phi-encoding}) are differences over the same time interval, so $\Delta t$ cancels exactly. These two classical scalars, the flux amplitude and its rate of change, are then mapped into the quantum circuit via the encoding $\mathcal{M}$.


\section{Dense Embedding: Amplitude and Phase Encoding}
\label{sec:dense-embedding}

Physical observables are mapped to the full Bloch sphere coordinates $(\theta, \phi)$. Any single-qubit state is a point on the Bloch sphere~\cite{nielsen2010}, parameterised by a polar angle $\theta \in [0, \pi]$ and an azimuthal angle $\phi \in [0, 2\pi)$.

The polar angle $\theta$ is derived from a multi-indicator classical score, not from the stress flux $\Phi$ directly. This ensures each qubit starts with the exact same initial probability as the classical AMR detector, so that the QAOA refines this baseline using entanglement rather than replacing it. The stress flux $\Phi$ is used for the phase encoding $\phi$ and for the Hamiltonian coefficients (Chapter~\ref{chap:hamiltonian}), but not for $\theta$.

\subsection{Classical Multi-Indicator Score}
\label{sec:classical-score}

The amplitude encoding uses the same multi-indicator detection criteria as production MHD codes (FLASH~\cite{fryxell2000}, Athena~\cite{stone2008}, PLUTO~\cite{mignone2007}). Four MHD-AMR indicators are combined:
\begin{enumerate}
    \item \textbf{Vorticity:} $|\omega_z| = |\partial_x v_y - \partial_y v_x|$ \quad (shear layers, vortex cores, KH rolls)
    \item \textbf{Velocity divergence:} $|\nabla \cdot \mathbf{v}| = |\partial_x v_x + \partial_y v_y|$ \quad (compression, shocks)
    \item \textbf{Current density:} $|J_z|$ \quad (current sheets, reconnection sites)
    \item \textbf{L\"ohner estimator on $|\mathbf{B}|$:} A scale-free second-derivative error estimator~\cite{lohner1987}:
    \begin{equation}
        E_{\text{L\"oh}}(f) = \frac{\sqrt{(f''_{xx})^2 + (f''_{yy})^2}}{|f'_x| + |f'_y| + \epsilon_{\text{L\"oh}} |f| + \epsilon}
    \end{equation}
    which peaks at discontinuities and is insensitive to smooth gradients.
\end{enumerate}
Each indicator is normalised to $[0, 1]$ by its domain-wide maximum. The combined score is the root-mean-square (RMS) of the four normalised indicators:
\begin{equation}
    s(x,y) = \sqrt{\frac{s_{\text{vort}}^2 + s_{\text{div}}^2 + s_{J_z}^2 + s_{\text{L\"oh}}^2}{4}}
    \label{eq:classical-score}
\end{equation}
RMS is used instead of $\max$ to avoid over-refinement when a single noisy indicator fires: a cell needs multiple indicators to agree before the combined score triggers refinement.

\subsection{Amplitude Encoding ($\theta$)}

The polar angle $\theta$ is derived from the classical score $s$ such that the initial probability of measuring $\ket{1}$ equals the score exactly:
\begin{equation}
    \theta_{ij} = 2 \arcsin\!\left(\sqrt{s_{ij}}\right) \quad \Longrightarrow \quad P(\ket{1}) = \sin^2\!\left(\frac{\theta_{ij}}{2}\right) = s_{ij}
    \label{eq:theta-encoding}
\end{equation}
At initialisation (before any QAOA layers), each qubit has exactly the same probability of being flagged as the classical AMR detector. The QAOA refines this starting point using entanglement (plaquette and gradient topology in the Hamiltonian).

\subsection{Phase Encoding ($\phi$)}

The azimuthal phase $\phi$ encodes the temporal growth rate of the stress flux, enabling the VQA to detect constructive interference between adjacent cells whose instabilities grow simultaneously:
\begin{equation}
    \phi_{ij}(t) = \frac{\pi}{2} \cdot \tanh\!\left( \beta \frac{\Phi_{ij}(t) - \Phi_{ij}(t - \Delta t_{\text{hybrid}})}{\langle|\Delta\Phi|\rangle} \right)
    \label{eq:phi-encoding}
\end{equation}
Here, $\beta$ is a trainable hyperparameter controlling the sensitivity to rapid changes, and $\langle|\Delta\Phi|\rangle = \text{mean}(|\Phi(t) - \Phi(t-\Delta t_{\text{hybrid}})|)$ is the global average flux change, providing a dimensionless normalisation. This maps $\phi$ to the range $[-\frac{\pi}{2}, +\frac{\pi}{2}]$: $\phi \approx +\frac{\pi}{2}$ for locally growing instabilities, $\phi \approx 0$ for background evolution, and $\phi \approx -\frac{\pi}{2}$ for damping.

The quantum state of each edge qubit $\ket{q_{ij}}$ is initialised as:
\begin{equation}
    \ket{q_{ij}} = \cos\!\left(\frac{\theta_{ij}}{2}\right)\ket{0} + e^{i\phi_{ij}}\sin\!\left(\frac{\theta_{ij}}{2}\right)\ket{1}
\label{eq_init}
\end{equation}

Stable regions ($s \approx 0$) map to $\ket{0}$ (north pole), unstable regions ($s \approx 1$) to $\ket{1}$ (south pole), and the azimuthal phase captures the dynamical tendency.

\subsection{Gate-Level Implementation}
\label{sec:gate-level}

In the Qiskit implementation, the state~(\ref{eq_init}) is prepared using the general single-qubit rotation gate $R(\theta, \varphi)$, defined by:
\begin{equation}
    R(\theta, \varphi) = \cos\!\left(\frac{\theta}{2}\right) I - i \sin\!\left(\frac{\theta}{2}\right) \bigl(\cos(\varphi)\, X + \sin(\varphi)\, Y\bigr)
\end{equation}
Applying $R(\theta, \varphi)$ to $\ket{0}$ yields:
\begin{equation}
    R(\theta, \varphi)\ket{0} = \cos\!\left(\frac{\theta}{2}\right)\ket{0} + e^{i(\varphi - \pi/2)} \sin\!\left(\frac{\theta}{2}\right)\ket{1}
\end{equation}
To match the target state~(\ref{eq_init}) with azimuthal angle $\phi$, a phase offset of $\pi/2$ is applied: each qubit is initialised with the gate call $R(\theta_{ij}, \phi_{ij} + \pi/2)$, which ensures that the resulting state has the desired phase $e^{i\phi_{ij}}$ on the $\ket{1}$ amplitude. This convention is consistent across all edge qubits (both horizontal and vertical).


\chapter{The Hamiltonian of Instability}
\label{chap:hamiltonian}

The Hamiltonian ground state identifies plasma regions that are topologically stressed or approaching instability. The key assumption is that high-Reynolds plasma dynamics arise from a finite number of physically distinct anomaly types: shear flows, vortices, magnetic reconnection, and helical kink modes. Each anomaly has a characteristic spatial signature encoded as a term in the quantum Hamiltonian.

Classical AMR detects instabilities only after they produce large gradients. Q-HAS encodes the structure of the plasma state at time $t$ to predict which regions will require refinement.

\section{Architecture Overview: Decoupled $f \times g \times T_{rcf}$}
\label{sec:hamilt-architecture}

A key design principle of the Hamiltonian is the decoupled multiplicative architecture. Each coefficient is a fully multiplicative product of three independent factors:
\begin{equation}
    \text{Coeff} = \text{Weight} \times g(\text{topology}) \times f(\text{scale}) \times T_{rcf}(\text{signal})
\end{equation}
This enforces strict separation between:
\begin{itemize}
    \item \textbf{Magnitude} of the anomaly ($T_{rcf}$: threshold-relative contrast filter),
    \item \textbf{Global thermodynamic scaling} ($f$: normal--critical gate),
    \item \textbf{Local topology} ($g$: leaky sigmoid gates).
\end{itemize}

The $ZZ$ and $ZZZZ$ interaction weights are fixed at 2 and 1 respectively (not trainable). The single-qubit $Z$ bias uses an adaptive weight $\alpha_z = w_{z,\text{frac}} \times \text{median}(|C|, |K|)$, where $w_{z,\text{frac}}$ is a trainable parameter. This keeps the $Z$ term subordinate to the spatial correlations ($ZZ$, $ZZZZ$) while breaking the QAOA ground-state degeneracy (without $Z$, the ferromagnetic $ZZ$/$ZZZZ$ coupling makes all-$\ket{0}$ and all-$\ket{1}$ degenerate).

The trainable parameters are: the encoding parameters ($\beta$), the uncertainty width ($\sigma$), the per-term contrast sensitivities ($\beta_{\text{curl}}$, $\beta_{\text{xpoint}}$), and the $Z$-bias fraction ($w_{z,\text{frac}}$), giving 5 trainable hyperparameters in total. Several parameters are frozen from prior calibration: $\text{threshold}_{\text{amr}} = 0.1496$ (from classical training), $\gamma_{\text{hydro}} = 2.0$, $\gamma_{\text{mag}} = 0.5$, and $\kappa = 10.0$. The $ZZ$ gradient coupling uses a fixed sensitivity $\beta = 1.0$ (its effective activation is controlled by $\sigma$ instead). All normalisation uses absolute physical non-dimensionalisation (dividing by fixed physical constants such as $Re_{\text{crit}} = 1$, $Rm_{\text{crit}} = 1$), never relative domain normalisation ($x / \max(x)$), ensuring that the Hamiltonian coefficients are stable across different grid resolutions and simulation times.

\section{The Hamiltonian Formula}

The Hamiltonian $H$ encodes the energy landscape of the plasma's topological stress. The VQA minimises $H$ to identify the configuration of stable and unstable regions most consistent with MHD anomaly physics.

QAOA~\cite{farhi2014} is used because: (1) the flux-graph formulation maps instability detection onto combinatorial graph optimisation (QAOA's native problem class), and (2) the Hamiltonian structure resembles the Ising model~\cite{lucas2014} (quantum Ising Hamiltonian are optimal computationally for QAOA because there is no need to diagonalise the Hamiltonian). QAOA alternates between two unitary operators:
\begin{itemize}
    \item The \textbf{structural (cost) Hamiltonian} $H_{\text{struct}}$ imprints a phase on each basis state proportional to its energy, marking configurations that correspond to physical anomalies.
    \item The \textbf{mixer Hamiltonian} $H_{\text{mixer}}$ generates quantum fluctuations (basis state transitions). Interference amplifies low-energy states and suppresses high-energy states.
\end{itemize}
Over $p$ layers, QAOA converges toward the ground state of $H_{\text{struct}}$.


\begin{equation}
    H_{\text{struct}} =
    \underbrace{\sum_i h_i Z_i}_{\substack{\text{Activity Bias} \\ \text{(adaptive } \alpha_z \text{)}}} +
    \underbrace{\sum_{\langle i,j \rangle} C_{ij} Z_i Z_j}_{\substack{\text{Gradient Coupling} \\ \text{(weight 2)}}} +
    \underbrace{\sum_{p} K_p \prod_{l \in p} Z_l}_{\substack{\text{Circulation Plaquette} \\ \text{(weight 1)}}} +
    \underbrace{\sum_{p} K_{\text{xpoint},p} \prod_{l \in p} Z_l}_{\substack{\text{X-point Reconnection} \\ \text{(weight 1, optional)}}}
\end{equation}

\noindent The X-point reconnection term detects hyperbolic magnetic null points via the determinant of the magnetic Jacobian (see Section~\ref{sec:xpoint-detection} below). It is a 4-body plaquette operator enabled by the \texttt{AdvAnomaliesEnable} flag. A fifth term for kink modes (Dzyaloshinskii--Moriya $XY - YX$ interactions driven by current helicity $\mathbf{J} \cdot \mathbf{B}$) is defined in the theoretical formulation but disabled in the 2D implementation (because kink modes are inherently three-dimensional, requiring helical deformation along the toroidal axis). The term is retained for future extension to 3D geometries. Note that $XY - YX$ is not diagonalizable in the $Z$ basis optimal for the QAOA algorithm (it is thus less optimal computationally to implement and run).

\noindent In this $H_{\text{struct}}$, each term acts as a specific ``sensor'' for a class of plasma instabilities, detailed below.

\begin{equation}
    H_{\text{mixer}} = \prod_i X_i
\end{equation}



\section{Activity Bias Term (The Validity Sensor)}
This single-body term biases each qubit according to whether its classical instability score exceeds the AMR threshold.
\begin{equation}
    H_{\text{bias}} = \sum_{i} h_i Z_i
\end{equation}
\textbf{Physical Coefficient $h_i$:} The bias coefficient directly encodes the deviation of the classical multi-indicator score (Equation~\ref{eq:classical-score}) from the AMR recursion threshold:
\begin{equation}
    h_i = \alpha_z \cdot \bigl(s_i - \text{threshold}_{\text{amr}}\bigr), \qquad \alpha_z = w_{z,\text{frac}} \times \text{median}\bigl(|C|, |K|\bigr)
\end{equation}
where $w_{z,\text{frac}}$ is a trainable hyperparameter and $\text{threshold}_{\text{amr}}$ is a Tier~1 hyperparameter optimised during training. The adaptive weight $\alpha_z$ scales with the median magnitude of the $ZZ$ and $ZZZZ$ coefficients, ensuring that the $Z$ term remains subordinate to the spatial correlations. Without the $Z$ term, the ferromagnetic $ZZ$/$ZZZZ$ coupling produces a degenerate ground state (all-$\ket{0}$ = all-$\ket{1}$); the adaptive bias breaks this degeneracy and forbids the QAOA from flipping every qubit to the same state.
\begin{itemize}
    \item When $s_i > \text{threshold}_{\text{amr}}$ (high instability score), $h_i > 0$: the qubit is biased toward $\ket{1}$ (flagged as active, rewards refinement).
    \item When $s_i < \text{threshold}_{\text{amr}}$ (low score), $h_i < 0$: the qubit is biased toward $\ket{0}$ (flagged as calm, penalises unnecessary refinement).
\end{itemize}
Since $\theta$ is initialised from the same score $s$ (Equation~\ref{eq:theta-encoding}), the Z-bias and the initial qubit amplitude are \textit{aligned}: classically active regions start closer to $\ket{1}$ and receive a positive Hamiltonian bias. The QAOA refines this starting point via the entanglement-based interaction terms. Basically, it only allows the other operators to flip the qubits close enough to the threshold where uncertainty lies.

\section{Gradient Coupling Term (The Boundary Sensor)}
This 2-body interaction term detects spatial discontinuities between neighbouring edges. A strong gradient indicates a boundary (shear layer, shock front, vortex boundary) that requires refinement.
\begin{equation}
    H_{\text{grad}} = \sum_{\langle ij, kl \rangle} C_{ij} Z_{ij} Z_{kl}
\end{equation}
\textbf{Physical Coupling $C_{ij}$:} The coupling coefficient follows the decoupled $f \times g \times T_{rcf}$ architecture, additionally modulated by an uncertainty-weighted Gaussian centred at the AMR decision boundary. Each coefficient is a fully multiplicative product that enforces strict separation between topology, global scaling, and local signal contrast:
\begin{equation}
    C_{ij} = 2 \cdot g_{\text{strain}}(Q_{\text{OW}}) \cdot \left\|\begin{pmatrix} f_{\text{hydro}}(Re) \cdot T_{\text{rcf}}(\Delta v) \\ f_{\text{mag}}(Rm) \cdot T_{\text{rcf}}(\Delta B) \end{pmatrix}\right\|
    \cdot \exp\!\left(-\left(\frac{\bar{s}_{ij} - \text{threshold}_{\text{amr}}}{\sigma}\right)^{\!2}\right)
\end{equation}
where $\bar{s}_{ij}$ is the average classical score of the two cells connected by the edge, the weight 2 is fixed (not trainable), and the $ZZ$ gradient terms use a fixed sensitivity $\beta = 1.0$ in their $T_{rcf}$ factor (sensitivity is controlled by $\sigma$ instead):
\begin{itemize}
    \item The uncertainty-weighted Gaussian $\exp\bigl(-\bigl((\bar{s} - \text{threshold}_{\text{amr}})/\sigma\bigr)^2\bigr)$ concentrates ZZ coupling near the decision boundary where the classical score is genuinely uncertain ($\bar{s} \approx \text{threshold}_{\text{amr}}$). Far from the threshold, where the classical decision is already confident, ZZ coupling is exponentially suppressed and the Z-bias alone drives decisions. This prevents ZZ from redundantly enforcing agreement where the classical detector already knows the answer, and focuses quantum corrections on the genuinely ambiguous cells. The trainable parameter $\sigma$ (range $[0.01, 0.30]$) controls the width of this activation band.
\end{itemize}
\begin{itemize}
    \item $g_{\text{strain}}(Q_{\text{OW}})$ is a leaky sigmoid topological gate based on the Okubo--Weiss Q-criterion~\cite{okubo1970, weiss1991} ($Q_{\text{OW}} > 0$ identifies strain-dominated regions):
    \begin{equation}
        g_{\text{strain}}(Q) = \frac{1}{1 + e^{-\kappa \, Q / Q_{\text{crit}}}}
    \end{equation}
    with $\kappa$ (trainable) controlling the steepness of the transition.
    \item $f(x)$ is a Normal-Critical gate providing logarithmic growth above the turbulence threshold:
    \begin{equation}
        f(x, x_{\text{crit}}, \gamma) = \begin{cases} x / x_{\text{crit}} & x \leq x_{\text{crit}} \\ 1 + \gamma \ln(x / x_{\text{crit}}) & x > x_{\text{crit}} \end{cases}
    \end{equation}
    The function is continuous at $x = x_{\text{crit}}$ (both sides evaluate to 1.0). Applied separately for the hydrodynamic ($f_{\text{hydro}}(Re, Re_{\text{crit}}, \gamma_{\text{hydro}})$) and magnetic ($f_{\text{mag}}(Rm, Rm_{\text{crit}}, \gamma_{\text{mag}})$) channels, with fixed critical thresholds $Re_{\text{crit}} = 1$ and $Rm_{\text{crit}} = 1$. The logarithmic form bounds growth: $Re = 3000$, $Re_{\text{crit}} = 1$, $\gamma = 2$ gives $f \approx 17$ rather than $\infty$, preventing any single edge from dominating the Hamiltonian.
    \item $T_{rcf}(\text{val}, \text{val}_{\text{crit}}, \beta)$ is a threshold-relative contrast filter that compares each signal to a fixed physical critical value:
    \begin{equation}
        T_{rcf}(\text{val}, \text{val}_{\text{crit}}, \beta) = \beta \cdot \max\!\left(0,\; \frac{\text{val}}{\text{val}_{\text{crit}}} - 1\right)
    \end{equation}
    Unlike the Michelson contrast (which compared $\text{val}$ to its spatial average, killing the signal when the domain is uniformly active), $T_{rcf}$ compares to a \textit{fixed physical threshold} ($\text{val}_{\text{crit}} = Re_{\text{crit}} \cdot \nu / dx$ for the gradient term). The signal survives in pre-anomaly regimes as long as $\text{val} > \text{val}_{\text{crit}}$, regardless of spatial uniformity. For the gradient ($ZZ$) term, the sensitivity $\beta = 1.0$ is fixed; the effective activation is controlled by the uncertainty width $\sigma$ instead. For the plaquette ($ZZZZ$) terms, the sensitivities $\beta_{\text{curl}}$ and $\beta_{\text{xpoint}}$ are trainable and specific to each term.
\end{itemize}
The four-factor decomposition ($f \times g \times T_{rcf} \times \text{Gaussian}$) enforces strict separation of concerns:
\begin{itemize}
    \item Below the physical threshold ($\text{val} < \text{val}_{\text{crit}}$), $T_{rcf} = 0$ and $C_{ij} = 0$: no constraint is imposed.
    \item In rotation-dominated regions ($Q_{\text{OW}} < 0$), the strain gate suppresses $C_{ij}$ to $0$, preventing false gradient triggers inside vortex cores while maintaining differentiability.
    \item At sharp strain boundaries ($Q_{\text{OW}} > 0$, high $Re$ or $Rm$, signal above threshold), all three factors are simultaneously large, producing strong ferromagnetic coupling that enforces neighbour agreement on refinement decisions.
    \item Far from the decision boundary ($|\bar{s} - \text{threshold}_{\text{amr}}| \gg \sigma$), the Gaussian factor suppresses $C_{ij}$ exponentially, so the classical score alone drives decisions for clear-cut cells. Near the boundary ($\bar{s} \approx \text{threshold}_{\text{amr}}$), the Gaussian factor is $\approx 1$ and ZZ correlations are fully active.
\end{itemize}

\section{Circulation Plaquette Term (The Vortex and Current Sensor)}
Vorticity and current density are detected via 4-body interactions on closed loops (plaquettes), implementing a discrete Stokes theorem: the circulation $\oint \mathbf{F} \cdot d\mathbf{l}$ measures the enclosed curl.
\begin{equation}
    H_{\text{topology}} = \sum_{p \in \text{Plaquettes}} K_p \prod_{(ij) \in p} Z_{ij}
\end{equation}
\textbf{Physical Circulation via Tangential Projection:} Horizontal edges carry $v_x$ and $B_x$; vertical edges carry $v_y$ and $B_y$. This gives two independent physical curls:
\begin{align}
    \omega_z(i,j) &= v_x(i,j) - v_x(i{+}1,j) + v_y(i,j{+}1) - v_y(i,j) \approx (\partial_x v_y - \partial_y v_x) \cdot \Delta A \\
    J_{z,\text{curl}}(i,j) &= B_x(i,j) - B_x(i{+}1,j) + B_y(i,j{+}1) - B_y(i,j) \approx (\partial_x B_y - \partial_y B_x) \cdot \Delta A
\end{align}
where $\omega_z$ is the discrete vorticity, $J_{z,\text{curl}}$ is the discrete current density (Amp\`ere's law) and $\Delta A$ the "infinitesimal" area where the curl is evaluated. The plaquette coefficient follows the same decoupled $f \times g$ architecture as the gradient coupling, with independent topological gates for the fluid and magnetic components:
\begin{equation}
    K_p = \left\|\begin{pmatrix} g_{\text{rot}}(Q_{\text{OW}}) \cdot f_{\text{hydro}}(Re) \cdot T_{\text{rcf}}(\omega_z, \omega_{z,\text{crit}}, \beta_{\text{curl}}) \\ g_{\text{mag}}(|J_z|) \cdot f_{\text{mag}}(Rm) \cdot T_{\text{rcf}}(J_{z,\text{curl}}, J_{z,\text{crit}}, \beta_{\text{curl}}) \end{pmatrix}\right\|
\end{equation}
where the critical values are $\omega_{z,\text{crit}} = Re_{\text{crit}} \cdot \nu / dx^2$ and $J_{z,\text{crit}} = Rm_{\text{crit}} \cdot \eta / dx^2$, and:
\begin{itemize}
    \item $g_{\text{rot}}(Q_{\text{OW}})$ is a leaky sigmoid gate activated in rotation-dominated regions ($Q_{\text{OW}} < 0$), complementary to $g_{\text{strain}}$:
    \begin{equation}
        g_{\text{rot}}(Q) = \frac{1}{1 + e^{+\kappa \, Q / Q_{\text{crit}}}}
    \end{equation}
    \item $g_{\text{mag}}(|J_z|)$ is a leaky sigmoid gate activated when the current density exceeds a critical threshold:
    \begin{equation}
        g_{\text{mag}}(|J_z|) = \frac{1}{1 + e^{-\kappa \, (|J_z| / J_{\text{crit}} - 1)}}
    \end{equation}
    \item The $f$-gates and reference scales $v_0 = \langle |\mathbf{v}| \rangle$, $B_0 = \langle |\mathbf{B}| \rangle$ are as defined in the gradient coupling term.
\end{itemize}
This decoupling ensures that the fluid vorticity component is gated by $Q_{\text{OW}} < 0$ (rotation), while the magnetic current-sheet component is gated independently by $|J_z|$. A vortex without a current sheet activates only the $g_{\text{rot}}$ branch, and a current sheet without vorticity activates only the $g_{\text{mag}}$ branch.
\begin{itemize}
    \item $\omega_z$ detects vortices: regions of concentrated fluid rotation.
    \item $J_{z,\text{curl}}$ detects current sheets and magnetic reconnection sites.
    \item In uniform or irrotational regions, $K_p = 0$: no plaquette penalty is applied.
\end{itemize}

\section{X-point Reconnection \& Chirality Terms (Advanced Anomalies)}
\label{sec:xpoint-detection}
Additional Hamiltonian terms detect topological features of the magnetic field beyond standard turbulence.


\textbf{X-point Reconnection Detection (Plaquette Term, Optional):}
To detect magnetic reconnection at X-points (hyperbolic null points of the magnetic field), the Hamiltonian uses a 4-body plaquette operator based on the determinant of the magnetic Jacobian:
\begin{equation}
    H_{\text{xpoint}} = \sum_{p \in \text{Plaquettes}} K_{\text{xpoint},p} \prod_{l \in p} Z_l
\end{equation}
The coefficient is:
\begin{equation}
    K_{\text{xpoint}} = f_{\text{mag}}(Rm) \cdot T_{rcf}\!\bigl(\max(0,\; -\det(\nabla \mathbf{B})),\; \text{val}_{\text{crit}},\; \beta_{\text{xpoint}}\bigr)
\end{equation}
where $\det(\nabla \mathbf{B}) = (\partial_x B_x)(\partial_y B_y) - (\partial_x B_y)(\partial_y B_x)$ is the determinant of the magnetic field Jacobian. At an X-point, field lines form a hyperbolic pattern and $\det(\nabla \mathbf{B}) < 0$. The $\max(0, -\det)$ factor retains only topologically hyperbolic regions. The critical value for the threshold-contrast is $\text{val}_{\text{crit}} = (Rm_{\text{crit}} \cdot \eta / (dx \cdot B_0))^2$, and $\beta_{\text{xpoint}}$ is a trainable sensitivity specific to this term. The term is self-limiting: it activates only where the magnetic topology is genuinely hyperbolic, without requiring a separate topological gate. It is enabled by the \texttt{AdvAnomaliesEnable} flag and trained on the Island Coalescence scenario. The Jacobian determinant therefore acts as an intrinsic geometric filter that suppresses inactive magnetic islands and responds only to genuine reconnection sites, obviating the need for an additional heuristic gate.


\section{Summary of Anomaly Mapping}
Table~\ref{tab:mapping} summarises how classical MHD phenomena are mapped to quantum operators in Q-HAS.

\begin{table}[h!]
\centering
\begin{tabular}{|l|l|l|l|c|}
\hline
\textbf{Anomaly} & \textbf{Physical Origin} & \textbf{Term} & \textbf{Operator} & \textbf{Weight} \\ \hline
Activity Bias & $\alpha_z \cdot (s_i - \text{threshold}_{\text{amr}})$ & Z bias & $Z_i$ & adaptive $\alpha_z$ \\ \hline
Shear / Gradient & $g_{\text{strain}} \times f \times T_{rcf} \times \text{Gauss}(\sigma)$ & ZZ coupling & $Z_i Z_j$ & 2 (fixed) \\ \hline
Vortex / Current & $\omega_z$, $J_{z,\text{curl}}$ (Stokes) & ZZZZ plaquette & $\prod_{\text{loop}} Z_k$ & 1 (fixed) \\ \hline
X-point Reconnection & $f_{Rm} \times T_{rcf}(-\det(\nabla\mathbf{B}))$ & ZZZZ plaquette & $\prod_{\text{loop}} Z_k$ & 1 (optional) \\ \hline
Kink (3D only) & $\mathbf{J} \cdot \mathbf{B}$ & DM interaction & $X_i Y_j - Y_i X_j$ & disabled \\ \hline
\end{tabular}
\caption{The Q-HAS Dictionary: Mapping MHD Instabilities to Quantum Operators. The $ZZ$ and $ZZZZ$ weights are fixed by operator degree. The $Z$ bias uses an adaptive weight subordinate to the spatial correlations.}
\label{tab:mapping}
\end{table}

\section{The Variational Ansatz}
The complete quantum circuit consists of two distinct stages: a data-embedding step followed by the QAOA evolution.

\begin{enumerate}
    \item \textbf{Data Embedding (Warm Start):}
    Each qubit is initialised from the physical flux data (Eq.~\ref{eq_init}), not from a uniform superposition. This prepares a product state $\ket{\psi_0}$ approximating the classical flux distribution:
    \begin{equation}
        \ket{\psi_0} \equiv \ket{\psi_0(\vec{\theta},\vec{\phi})} = \bigotimes_{\langle i,j \rangle} \left( \cos\!\left(\frac{\theta_{ij}}{2}\right)\ket{0} + e^{i\phi_{ij}}\sin\!\left(\frac{\theta_{ij}}{2}\right)\ket{1} \right)
    \end{equation}
    This warm-start strategy constrains VQA exploration to the neighbourhood of the physical solution, reducing QAOA iterations and mitigating the barren plateau problem~\cite{mcclean2018}.

    \item \textbf{QAOA Evolution:}
    The circuit applies $p$ layers of alternating unitary evolutions governed by $H_{\text{struct}}$ and $H_{\text{mixer}}$. For depth $p$, the trial state is:
    \begin{equation}
        |\psi(\vec{\Omega},\vec{\Gamma})\rangle = \prod_{k=1}^{p} \left( U_{\text{mixer}}(\Omega_k)\, U_{\text{struct}}(\Gamma_k) \right) |\psi_0(\vec{\theta},\vec{\phi})\rangle
    \end{equation}

    where:
    \begin{itemize}
        \item \textbf{The Phase Separator:} $U_{\text{struct}}(\Gamma_k) = e^{-i\Gamma_k H_{\text{struct}}}$ embeds the MHD constraints (shear, topology, reconnection) into the quantum phase of each basis state.
        \item \textbf{The Mixer:} $U_{\text{mixer}}(\Omega_k) = e^{-i\Omega_k H_{\text{mixer}}} = \prod_i e^{-i\Omega_k X_i}$ generates quantum fluctuations (spin flips) and facilitates interference between basis states, guided by the initial phases $\phi$.
    \end{itemize}

The variational parameters $\vec{\Gamma}, \vec{\Omega}$ are optimised by a classical loop (using the COBYLA optimiser~\cite{powell1994}) to minimise the energy expectation value $E = \langle \psi | H_{\text{struct}} | \psi \rangle$. The resulting ground state identifies the configuration of stable and unstable qubits that best fits the encoded physics.

\end{enumerate}



\section{Typical quantum circuit}

\begin{figure}[h]
\centering

\mbox{
    \Qcircuit @C=0.8em @R=1.2em {
        % ,  Line 1: Top Qubit , 
        \lstick{\ket{0}} 
        & \gate{R(\theta_{1,1,v},\phi_{1,1,v})} 
        & \multigate{3}{e^{i\Gamma_1H_{struct}}} % Couvre cette ligne + 3 en dessous (Total 4)
        & \multigate{3}{e^{i\Omega_1H_{mixer}}} 
        & \qw
        & \cdots 
        & {}
        & \qw
        & \multigate{3}{e^{i\Gamma_pH_{struct}}} 
        & \multigate{3}{e^{i\Omega_pH_{mixer}}} 
        & \qw
        & \meter \ar @{=>} [4,0]
        & \qw
        & \qw
        & \qw 
        \\
        % ,  Line 2: Middle Upper , 
        \lstick{\ket{0}} 
        & \gate{R(\theta_{1,1,h},\phi_{1,1,h})}  
        & \ghost{e^{i\Gamma_1H_{struct}}}
        & \ghost{e^{i\Omega_1H_{mixer}}}  
        & \qw
        & \cdots 
        & {}
        & \qw
        & \ghost{e^{i\Gamma_1H_{struct}}} 
        & \ghost{e^{i\Omega_pH_{mixer}}} 
        & \qw 
        & \qw
        & \meter \ar @{=>} [3,0]
        & \qw
        & \qw
        \\
        % ,  Line 3: Middle Lower , 
        \lstick{\vdots\hspace{0.2 cm}}
        & \vdots \hspace{0.1 cm}
        & {}
        & {} 
        & {}
        & \vdots 
        & {}
        & {}
        & {} 
        & {}
        & {} 
        & {}
        & {}
        & \ddots
        & {}
        \\
        % ,  Line 4: Bottom Qubit , 
        \lstick{\ket{0}} 
        & \gate{R(\theta_{n,n,h},\phi_{n,n,h})} 
        & \ghost{e^{i\Gamma_1H_{struct}}} 
        & \ghost{e^{i\Omega_1H_{mixer}}} 
        & \qw
        & \cdots 
        & {}
        & \qw
        & \ghost{e^{i\Gamma_1H_{struct}}} 
        & \ghost{e^{i\Omega_pH_{mixer}}} 
        & \qw
        & \qw
        & \qw
        & \qw
        & \meter \ar @{=>} [1,0] % Descend de 1 ligne vers le Bus
        \\
        % ,  Line 5: Classical Bus , 
        \lstick{Bus}
        & \cw {/}
        & \cw 
        & \cw 
        & \cw
        & \cdots 
        & \dstick{\raisebox{-0.3em}{\hspace{-1em}$\underbrace{\hspace{27.5em}}_{p \text{ layers}}$}}
        & \cw
        & \cw 
        & \cw
        & \cw
        & \cw
        & \cw
        & \cw
        & \cw \\
    }
}
\vspace{0.7cm}
\caption{Minimization of $\bra{\psi(\vec{\Omega},\vec{\Gamma})}H_{struct}\ket{\psi(\vec{\Omega},\vec{\Gamma})}$ with the QAOA circuit with Warm Start}
\end{figure}


\chapter{The Algorithm Structure}
\label{chap:algorithm}

Current quantum hardware limits the number of available qubits, making full-grid QAOA infeasible. Q-HAS uses a hierarchical (multi-level) scanning strategy via breadth-first search (BFS) over the spatial domain.

\section{Level-by-Level BFS Scanning}

The number of qubits is fixed, determining the patch size: for a $\zeta$-dimensional periodic grid, $\zeta N_{\text{lim}}^{\zeta}$ qubits cover $N_{\text{lim}}^{\zeta}$ coarse blocks. All patches at the same depth are processed before moving to the next level, enabling batched execution.

At each depth level, physics is downsampled to the $N_{\text{lim}} \times N_{\text{lim}}$ grid and the VQA returns a probability map (averaging horizontal and vertical edge marginals). For each sub-block:
\begin{itemize}
    \item If $\text{prob} \geq \text{threshold}_{\text{eff}}$: the sub-block is subdivided and added to the next depth level.
    \item If $\text{prob} < \text{threshold}_{\text{eff}}$: the sub-block is marked as a coarse leaf (calm region, solved at coarse resolution).
\end{itemize}

Recursion terminates at maximum depth $D_{\text{max}}$ or when the patch size falls below $\text{min\_size} = 6$ pixels. The termination depth determines the local solver resolution. The threshold is constant over depth (attempted to increase it linearly with depth at first to prevent excessive fine-scale refinement but a constant value over depth showed better results).\\

\textbf{Decoupled scan/solve depths:} The scan depth (how many VQA levels are evaluated) can be capped independently from the solve depth (the resolution hierarchy used by the solver). Patches identified by the VQA are stored with their natural solve depth, ensuring that the solver computes the correct \texttt{local\_factor} even when the VQA scan stops early.

\section{Temporal Memory: Time-To-Live (TTL)}

An anomaly detected at time $t$ may weaken at $t + \Delta t_{\text{hybrid}}$ due to downsampling, noise, or transient fluctuation. Without temporal memory, the patch would be immediately dropped (false negative).

Each refined patch carries a \textbf{TTL (Time-To-Live) counter} (default $\tau = 1$ hybrid step). When flagged for refinement, TTL is set to $\tau$. At each subsequent step:
\begin{itemize}
    \item If the VQA probability exceeds the threshold: TTL is \textit{reset} to $\tau$ (fresh detection).
    \item If the VQA probability drops below the threshold but $\text{TTL} > 0$: the patch \textit{survives} and TTL is decremented by 1.
    \item If $\text{TTL} = 0$: the patch is released to the coarse level.
\end{itemize}
This stabilises the refinement tree: real instabilities have time to develop while transient false positives decay. The TTL map (keyed by patch bounds) persists across hybrid steps. Both VQA-guided and classical AMR paths maintain independent TTL maps.

\section{Boundary Probing}

After VQA execution, the engine checks whether edge qubits (patch boundary) have elevated activation probability. If an anomaly touches the boundary, the corresponding sub-cell is probed deeper even at marginal probability ($\geq 50\%$ of the effective threshold). This solves the ``lost topology when zooming'' problem: the algorithm rediscovers anomalies locally rather than propagating from the parent.

The boundary detector compares the mean probability of each edge row/column (top, bottom, left, right) against the interior mean plus a tolerance. Active boundaries force the adjacent sub-cell into the next refinement level.

\begin{algorithm}[H]
\caption{Q-HAS: Time-Stepping Loop}
\begin{algorithmic}[1]
\State \textbf{Initialize:} Classical Plasma State $\mathbf{U}_0$ on deep grid $G_0$ (resolution of $G_0$ is $N^{\zeta}$, with $\zeta$ the dimension of the simulation [2D or 3D])
\State Qubits number = $\zeta N_{lim}^{\zeta}$ 
\State $D_{max}= \lceil\frac{\log(N)}{log(N_{lim})}\rceil$
\State \textbf{Hyperparameters:} Convergence tolerance of the Quantum Approximate Optimization Algorithm (QAOA) $\epsilon$, Hamiltonian terms sensibilities and state sensibilities.
\For{$t = 0$ to $T_{end}$}
    \State \textit{// Step 1: Classical Coarse Update}
    
    \If{$t_{hybrid}|t$}
        \State \textit{// Step 1: Define coarse grid} 
        \State $G_{coarse}= \text{Reshape}(G_0, N_{lim})$
        \State \textit{// Step 2: Quantum Risk Assessment}
        \State Extract Fluxes $\Phi = F_{extract}(\mathbf{U}_{t},G_{coarse})$.
        \State Map Fluxes to Graph Parameters: $\mathbf{\theta, \phi} \leftarrow \mathcal{M}(\Phi)$.
        \State \textbf{VQA Routine:}
        \State \quad Minimize $\langle\psi(\vec{\Omega}, \vec{\Gamma})| H_{struct}|\psi(\vec{\Omega}, \vec{\Gamma})\rangle$ with QAOA.
        \State \quad Measure Qubits $\to$ Collapse to bitstring $z^*$.
        
        \State \textit{// Step 3: Identify High-Risk Zones}
        \State Let $\mathcal{R}$ be the set of edges where $z^*$ indicates instability (e.g., $|1\rangle$).
        
        \State \textit{// Step 4: Recursive Refinement (Zoom)}
        \If{$\mathcal{R}$ is not empty AND Depth $< D_{max}$}
            \State $\Omega_{sub} \leftarrow \text{Zoom}(G_0, \mathcal{R})$ \Comment{Define the spatial window in the high-resolution grid}
            \State $G_{coarse} \leftarrow \text{Reshape}(\Omega_{sub}, N_{lim})$ \Comment{Resample this area in $N_{lim}^d$ points}
            \State $\text{Depth} \leftarrow \text{Depth} + 1$
            \State \textbf{Recurse} Step 2 with the new grid $G_{coarse}$
        \EndIf
        \State Apply Fine Mesh AMR on final regions $G =\text{AMR}(G_0, \mathcal{R}) $.
    \EndIf
    
    \State \textit{// Step 5: Correction \& Advance}
    
    \State $\mathbf{U}_{t} \leftarrow \text{ClassicalSolver}(\mathbf{U}_t, G,  \Delta t_{coarse})$
\EndFor
\end{algorithmic}
\end{algorithm}

\chapter{Theoretical Framework and Optimisation Dynamics}
\label{chap:theory}
\section{Description of the Algorithm}

Q-HAS reduces computational cost by identifying precisely which regions require high-resolution computation, rather than solving the MHD equations faster. This differs from DNS (full-resolution everywhere) and classical AMR (reactive gradient-based refinement).

\paragraph{Localised solving.}
In most plasma simulations, only a small fraction of the domain generates anomalies. The rest can be evolved at coarse resolution with minimal accuracy loss. Q-HAS assigns resolution proportional to the detection depth, so both location and severity of instabilities inform the refinement decision.

\paragraph{Unstable zone identification.}
Detection proceeds in two stages:

\noindent \textit{Stage 1 (Naive detection):} The flux amplitudes $\Phi_{ij}$ identify regions with strong gradients, which are direct indicators of instabilities. A classical gradient-based AMR would stop here.

\noindent \textit{Stage 2 (Topological validation):} The Hamiltonian $H_{\text{struct}}$ evaluates whether the pattern of stable and unstable qubits is consistent with known MHD anomaly structures. If flipping a qubit from stable ($\ket{0}$) to unstable ($\ket{1}$) reduces the Hamiltonian energy, meaning the flip completes a recognised anomaly pattern, then the flip is energetically favoured. This contextual analysis filters out incoherent noise and identifies genuine topological structures.

Two aspects provide the quantum advantage:
\begin{itemize}
    \item \textbf{Quantum tunnelling:} The QAOA mixer Hamiltonian enables the optimiser to escape local minima through quantum tunnelling, exploring configurations that a greedy classical search would miss.
    \item \textbf{Phase-based anticipation:} If two adjacent regions have high growth rate ($\dot{\Phi}$) but low amplitude ($\Phi$), a classical gradient detector ignores them. In Q-HAS, the encoded phases $\phi$ align, creating a low-energy interference pattern captured by the mixer. This detects a pre-instability resonance: an anomaly forming but not yet at visible amplitude. The predictive advantage is identifying the dynamical tendency before the amplitude becomes critical.
\end{itemize}

The hierarchical structure provides an additional efficiency gain: determining the precise depth of each anomaly avoids unnecessary mesh refinement in surrounding regions.

\section{Initialization of Structural Parameters (of the QAOA)}
\label{sec:initialization}
\subsection{Initialization of $\vec{\Gamma}$}

While the Null-Gradient Proposition proves we cannot \textit{optimise} $\Gamma$ initially, we must still \textit{set} it. If $\Gamma=0$, the Mixer force $\nabla_\Omega$ also vanishes (as the state remains an eigenstate of $X$). Therefore, $\Gamma$ must be initialised to define the ``target phase'' that $\Omega$ will act upon.

We propose a Linear-Ramp ``Phase-Matching'' Strategy for general depth $p$.

\subsubsection{The Problem of Phase Wrapping}
The unitary operator for the cost function is $U_C = e^{-i \gamma H_{struct}}$.
The eigenvalues of $H_{struct}$ (the problem energies) map to phases $\varphi = \gamma E$.
\begin{itemize}
    \item \textbf{Risk:} If $\gamma$ is too large, $\gamma E > 2\pi$. High-energy states "wrap around" and become indistinguishable from low-energy states (aliasing).
    \item \textbf{Constraint:} To preserve the gradient signal, the total phase rotation across the entire circuit must stay within the principle branch $[0, \pi]$.
\end{itemize}

\subsubsection{Optimal Initialization Rule}
Let $E_{max}$ be the maximum expected eigenvalue of $H_{struct}$ (approximated by the sum of absolute weights of all edges).
The total structural angle $\Gamma_{total} = \sum_{k=1}^p \Gamma_k$ should satisfy:
\begin{equation}
    \Gamma_{total} \approx \frac{\pi}{E_{max}}
\end{equation}

\begin{equation}
    E_{\max} \approx \sum_i |h_i| + \sum_{\langle i,j \rangle} |C_{ij}| + \sum_p |K_p| + \sum_v |\Delta_v| + 2 \sum_{\langle i,j \rangle} |D_{ij}|
\end{equation}

For a circuit of depth $p$, we distribute this total angle using a Linear Ramp Schedule (Trotterized Adiabatic Evolution):
\begin{equation}
    \Gamma_k^{init} = \left( \frac{2k}{p(p+1)} \right) \cdot \Gamma_{total} \quad \text{for } k=1 \dots p
\end{equation}
\textit{(Note: The pre-factor normalizes the sum to 1).}

\textbf{Physical Justification for the Ramp ($\Gamma_p < \Gamma_q \text{ if } p<q$) :}
\begin{enumerate}
    \item \textbf{Layer 1 (Small $\Gamma$):} Acts as a linear approximation. It applies a small "tilt" to the landscape, creating a coherent gradient direction for the first Mixer to follow without scrambling the phases.
    \item \textbf{Layer $p$ (Large $\Gamma$):} Acts as a constraint enforcer. As the system evolves, we increase the weight of the constraints to "lock in" the solution, mimicking the adiabatic process where $H(t)$ transitions from $H_{mix}$ to $H_{struct}$.
\end{enumerate}

This initialization ensures that the initial $\Omega$ gradients are non-zero, coherent, and point towards the global minimum rather than a local trap.

\subsection{Initialization of $\vec{\Omega}$}

The minimisation of $H_{\text{struct}}$ via the QAOA ansatz selects quantum states corresponding to physical instabilities. The following derivations show how the energy gradient drives state evolution and why the Mixer ($\Omega$) is mandatory for algorithm startup.

\subsubsection{The Null-Gradient Proposition (Structural Parameters)}
\label{sec:null-gradient}
The warm-start initialisation prevents classical optimisation from adjusting the structural parameters $\Gamma$ first; it must activate quantum fluctuations via $\Omega$.

\textbf{Proposition:} At the initialization point where all mixer parameters are set to zero ($\vec{\Omega} = \vec{0}$), the gradient of the energy with respect to any structural parameter $\Gamma_k$ is strictly zero, regardless of the circuit depth $p$.

\textbf{Proof:}
Let the depth be $p \geq 1$. The ansatz is:
\begin{equation}
    \ket{\psi(\vec{\Omega}, \vec{\Gamma})} = \prod_{k=1}^p \left( e^{-i\Omega_k H_{mixer}} e^{-i\Gamma_k H_{struct}} \right) \ket{\psi_0}
\end{equation}
At initialization, we set $\vec{\Omega} = \vec{0}$, so $e^{-i(0) H_{mixer}} = \mathbb{I}$. The state simplifies to a product of commuting exponentials:
\begin{equation}
    \ket{\psi(\vec{0}, \vec{\Gamma})} = \left( \prod_{k=1}^p e^{-i\Gamma_k H_{struct}} \right) \ket{\psi_0} = e^{-i (\sum_k \Gamma_k) H_{struct}} \ket{\psi_0}
\end{equation}
The energy expectation value is $E\equiv E(\vec{\Omega},\vec{\Gamma}) = \bra{\psi} H_{struct} \ket{\psi}$. Substituting the simplified state:
\begin{equation}
    E = \bra{\psi_0} e^{i (\sum \Gamma) H_{struct}} H_{struct} e^{-i (\sum \Gamma) H_{struct}} \ket{\psi_0}
\end{equation}
Since any operator commutes with itself ($[H_{struct}, e^{-i \alpha H_{struct}}] = 0$), the unitary operators cancel out:
\begin{equation}
    E = \bra{\psi_0} H_{struct} \ket{\psi_0} = \text{Constant}
\end{equation}
Consequently, the gradient is null:
\begin{equation}
    \nabla_{\Gamma_k} E \Big|_{\vec{\Omega}=0} = \frac{\partial}{\partial \Gamma_k} (\text{Const}) = 0 \quad \forall k \in [1, p]
\end{equation}
\textbf{Physical Implication:} The algorithm cannot lower energy by adjusting constraint weights alone ($\Omega = 0$). It must activate the Mixer ($\Omega$) to break symmetry. Phase-driven prediction is the exclusive engine of algorithm startup.

\section{Optimisation Dynamics: Gradient Structure and Symmetry Breaking}
\label{sec:minimization-dynamics}

Throughout this section we analyse the variational landscape of the QAOA energy $E(\vec\Omega,\vec\Gamma) = \langle\psi(\vec\Omega,\vec\Gamma)|H_{\text{struct}}|\psi(\vec\Omega,\vec\Gamma)\rangle$. We emphasise that the results below are \emph{structural} statements about the gradient: they describe where the landscape is flat, in which direction it first tilts, and what physical quantities drive the first-order change. They do \emph{not} constitute a proof of computational speedup, query-complexity separation, or convergence to the ground state, and they should not be read as establishing a quantum advantage. Whether the framework realises a quantitative improvement over the classical multi-indicator baseline is an empirical question addressed in Chapters~\ref{chap:results} and~\ref{chap:falsification}.

\subsection{Dynamics when $\vec{\Omega}= 0$ (initialization value)}
\subsubsection{The Driving Force (General Mixer Gradient)}
Since $\nabla_\Gamma E = 0$, optimisation is entirely driven by the mixer parameters $\vec{\Omega}$. The following generalises the gradient for any layer $k \in [1, p]$ at $\vec{\Omega} = \vec{0}$.

\textbf{Proposition (Mixer-gradient structure):} The gradient $\partial_{\Omega_k} E$ is proportional to the expectation value of the commutator $[H_{\text{struct}}, H_{\text{mixer}}]$ on the phase-shifted state up to layer $k$.

\textbf{Proof:}
The unitary evolution after layer $k$ is $U_{>k}$. At $\vec{\Omega}=\vec{0}$, this consists solely of structural gates:
\begin{equation}
    U_{>k} = \prod_{j=k+1}^p e^{-i\Gamma_j H_{struct}} = e^{-i (\sum_{j=k+1}^p \Gamma_j) H_{struct}}
\end{equation}
The effective Hamiltonian at layer $k$ is $H_{eff} = U_{>k}^\dagger H_{struct} U_{>k}$. Since $[H, e^{i\alpha H}] = 0$, future layers cancel:
\begin{equation}
    H_{eff} = H_{struct}
\end{equation}
Thus, the gradient depends only on the state arriving at layer $k$, denoted $\ket{\psi_{k-1}}$:
\begin{equation}
    \partial_{\Omega_k} E \big|_{\vec{\Omega}=0} = -i \bra{\psi_{k}} [H_{struct}, H_{mixer}] \ket{\psi_{k}}
\label{eq:mixer-gradient}
\end{equation}
where $\ket{\psi_{k}} = \left( \prod_{j=1}^{k} e^{-i\Gamma_j H_{struct}} \right) \ket{\psi_0}$.\\

\noindent\textbf{Physical Implication: Phase-Driven Ignition}\\
The structural operator $e^{-i\Gamma_k H_{struct}}$ acts as a \textit{phase oracle} before mixer activation.
\begin{itemize}
\item \textbf{Immediate Gradient Coupling ($k=1$):} $H_{struct}$ imprints topological defects (shear, vorticity) into relative phases. The commutator $[H_{struct}, H_{mixer}]$ converts this phase landscape into a non-zero gradient force. The physics of the problem dictates the initial optimisation direction.
\item \textbf{Recursive Phase Refinement ($k>1$):} The driving force is the accumulated phase distribution ($\phi_{\text{data}} + \phi_{\text{correction}}$). The mixer at layer $k$ exploits the interference pattern from all preceding structural gates to resolve residual symmetries.
\end{itemize}


\subsubsection{Detailed Derivation by Hamiltonian Term}
To understand exactly which physical phenomena trigger a quantum response, we compute the commutator $[H_{\text{term}}, X_k]$ for each component of $H_{\text{struct}}$. We use the Pauli identities $[Z, X] = 2iY$ and $[Y, X] = -2iZ$.

\paragraph{A. Shear Term (Viscosity)}
For a link between qubits $i$ and $j$: $H_{shear} = C_{ij} Z_i Z_j$. The shear energy gradient is the sum of the contribution of each qubit. The generic qubit $k$ contribution is: 
\begin{align}
    \nabla_{\Omega_m} E_{shear} &\propto -i \bra{\psi_m} [Z_i Z_j, X_i] \ket{\psi_m} \nonumber \\
    &= -i \bra{\psi_m} [Z_i, X_i] Z_j \ket{\psi_0} \nonumber \\
    &= -i (2i) \bra{\psi_m} Y_i Z_j \ket{\psi_m} = 2 C_{ij} \langle Y_i Z_j \rangle
\end{align}

\paragraph{B. Topology/X-point Term (Plaquette)}
For a generic multi-body term involving a set of qubits $\mathcal{S}$ (e.g., $ZZ$ or $ZZZZ$ operators): $H_{\mathcal{S}} = K \prod_{m \in \mathcal{S}} Z_m$.
If we calculate the associated energy gradient (related to the energy contained in the $\mathcal{S}$ topology) contribution term from qubit $k \in \mathcal{S}$:
\begin{equation}
    \nabla_{\Omega_m} E_{\mathcal{S}} \propto 2 K \left\langle \psi_m \left| Y_k \prod_{m \in \mathcal{S}, m \neq k} Z_m \right| \psi_m\right\rangle
\end{equation}

\subsubsection{From Gradient to Bit-Flip: The Physics of Detection}
\label{sec:coherence-filter}
How does a non-zero energy gradient transform a stable qubit ($\ket{0}$) into a flagged anomaly ($\ket{1}$)?

\paragraph{1. The Update Rule}
At optimisation step $t$:
\begin{equation}
    \Omega_k^{(t+1)} = \Omega_k^{(t)} - \eta \cdot \partial_{\Omega_k} E \big|_{\vec{\Omega}^{(t)}, \vec{\Gamma}^{(t)}}
\end{equation}
where $\eta$ is the learning rate.
Since we proved that $\nabla_{\Omega} E \neq 0$ at initialization (Equation~\ref{eq:mixer-gradient}), the first step of the algorithm is guaranteed to be non-trivial. If the gradient is negative (indicating an instability is detected), $\Omega_k$ becomes positive, activating the rotation $U_{mixer}$.

\paragraph{2. The Rotation (The Flip)}
The unitary operator is $U_{mixer}(\Omega) = e^{-i \Omega X}$. Geometrically, this is a rotation on the Bloch sphere around the X-axis.
\begin{itemize}
    \item \textbf{Start:} $\ket{\psi} \approx \ket{0}$ (North Pole).
    \item \textbf{Action:} A strong gradient induces a large $\Omega$. The operator rotates the vector towards the South Pole ($\ket{1}$).
    \item \textbf{Result:} The probability of measuring $\ket{1}$ becomes $P(1) = \sin^2(\Omega)$.
\end{itemize}

\begin{tcolorbox}
\paragraph{3. The Coherence Filter (Why context matters)}
The magnitude of the rotation $\Omega$ is directly proportional to the gradient force $\vec{F}_q$ on qubit $q$ derived in the following equation:
\begin{equation}
    \vec{F}_q \propto \sin(\theta_q)\cdot\underbrace{\sin(\phi_q)}_{\text{Phase Boost}} \cdot\underbrace{\prod_{m \neq q} \cos(\theta_m)}_{\text{Topological Coherence}}
\label{eq:gradient-force}
\end{equation}

\textbf{Scenario: Destructive Interference (Noise Suppression)}
Suppose qubit $k$ has a high phase $\phi_k$ (potential anomaly), but its neighbors are in a random, incoherent turbulent state.
\begin{itemize}
    \item The neighbors fluctuate randomly between $\ket{0}$ and $\ket{1}$, so $\cos(\theta_m)$ averages to 0.
    \item The product term $\prod \cos(\theta_m) \to 0$.
    \item \textbf{Consequence:} $\nabla_{\Omega} E \to 0$.
    \item \textbf{Mechanism:} The optimizer sees a zero gradient. $\Omega$ remains zero. No rotation is applied. The qubit remains at $\ket{0}$.
    \item \textbf{Physics:} The algorithm refuses to flip the bit because the anomaly is isolated and topologically inconsistent.
\end{itemize}

\textbf{Scenario: Constructive Interference (Structure Detection)}
Suppose qubit $k$ is part of a coherent vortex filament where neighbors are laminar ($\ket{0}, \cos \theta \approx 1$).
\begin{itemize}
    \item The coherence term $\prod \cos(\theta_m) \approx 1$.
    \item The gradient is maximal: $\nabla_{\Omega} E \propto \sin(\phi_k)$.
    \item \textbf{Mechanism:} The optimizer receives a strong signal. $\Omega$ increases significantly. The rotation $e^{-i \Omega X}$ is applied.
    \item \textbf{Physics:} The bit physically flips to $\ket{1}$, flagging the entire region for Adaptive Mesh Refinement.
\end{itemize}
\end{tcolorbox}
\subsection{Dynamics when $\vec{\Omega} \neq 0$ (Active Dynamics: The Competition of Torques)}

Once $\Omega_k > 0$, the null-gradient proposition for $\Gamma$ collapses. The system enters the active regime, governed by two competing unitary forces acting as orthogonal torques on the Bloch sphere.

\subsubsection{The Origin of Forces: Back-Action Mechanism}
The force on a qubit at layer $k$ arises from back-propagation of the final energy measurement through future layers. The effective Hamiltonian at layer $k$ is $H_{\text{struct}}$ rotated by all subsequent gates $U_{>k}$:
\begin{equation}
    H_{eff}^{(k)} = U_{>k}^\dagger H_{struct} U_{>k}
\end{equation}

The gradient with respect to any parameter $\Lambda_k$ ($\Gamma$ or $\Omega$) is:
\begin{equation}
    F_{\Lambda_k} = -\nabla_{\Lambda_k} E = -i \bra{\psi_{k}} [G_k, H_{eff}^{(k)}] \ket{\psi_{k}}
\end{equation}

\subsubsection{Force 1: The Interference Force ($\nabla_\Gamma$)}
This force acts on the structural parameter $\Gamma$.
\begin{itemize}
    \item \textbf{Location:} Acts locally on the Z-axis of the qubit (Phase rotation).
    \item \textbf{Origin:} It arises solely from the \textbf{Mixer's rotation of the basis} in future layers.
\end{itemize}

\textbf{Derivation:}
At $\Omega=0$, $H_{eff}$ commuted with $H_{struct}$, resulting in zero force. Now that $\Omega \neq 0$, the effective Hamiltonian contains terms non-diagonal in the Z-basis (e.g., $X$ and $Y$ terms).
\begin{equation}
    F_{\Gamma_k} = -i \bra{\psi_{k}} [H_{struct}, H_{eff}^{(k)}] \ket{\psi_{k}}
\end{equation}

\textbf{Physical Interpretation:}

$H_{struct}$ is diagonal in the computational basis with eigenvalues $E(z)$. The unitary $e^{-i\Gamma H_{struct}}$ imprints energy-dependent phases:
\begin{equation}
|z\rangle \rightarrow e^{-i\Gamma E(z)}|z\rangle .
\end{equation}

These phases become relevant when the mixer couples different basis states. Higher-energy configurations acquire larger phase shifts and interfere less coherently, suppressing their probability. Low-energy configurations remain coherent and accumulate probability weight. The force $F_\Gamma$ acts the following way: destructive interference suppresses high-energy configurations while constructive interference reinforces the low-energy sector.

\subsubsection{Force 2: The Tunnelling Force ($\nabla_\Omega$)}
This force acts on the mixer parameter $\Omega$.
\begin{itemize}
    \item \textbf{Location:} Acts locally on the X-axis of the qubit (Bit-flip rotation).
    \item \textbf{Origin:} It arises from the misalignment between the current quantum state vector and the principal axes of the effective Hamiltonian.
\end{itemize}

\textbf{Derivation:}
\begin{equation}
    F_{\Omega_k} = -i \bra{\psi_{k}} [H_{mixer}, H_{eff}^{(k)}] \ket{\psi_{k}}
\end{equation}

This commutator measures the virtual torque exerted by the problem constraints on the transverse axis.

$H_{mixer} \sim X$ couples basis states differing by a single bit flip. Classical configurations are local minima separated by energy barriers. The mixer-induced rotations spread the quantum state coherently across neighbouring configurations, enabling amplitude to appear beyond energy barriers without traversing high-energy intermediate states.

This is analogous to \textit{quantum tunnelling}: the state redistributes probability amplitude across otherwise inaccessible configurations. The magnitude of $F_\Omega$ determines exploration strength, allowing the algorithm to escape local minima.


\section{Algorithm Complexity Analysis}
\label{sec:complexity}

We compare the asymptotic cost of DNS, classical AMR and Q-HAS. Throughout, $N_{\text{fine}}$ denotes the linear resolution required to resolve the smallest dissipative scale and $d \in \{2,3\}$ the spatial dimension. We distinguish between \emph{cost per time step} and \emph{cost per unit simulated time}; both are relevant, but confusing them leads to the wrong exponents.

\paragraph{A note on what this section does and does not show.}
The bounds below are purely asymptotic and describe only how each algorithm's cost scales with $N_{\text{fine}}$ and with the sparsity factors $S_{\text{class}}, S_{\text{quant}}$. They do not establish a quantum advantage: (i) they assume that $T_{\text{VQA}}$ is effectively constant in the problem size, which hides the precision and depth dependence discussed below; (ii) they assume that $S_{\text{quant}} < S_{\text{class}}$ and that $\nu_{\text{scan}} < 1$ is sustainable, neither of which is verified empirically in Chapter~\ref{chap:results}. The section should be read as \emph{what it would take} for such an architecture to be cost-effective, not as a proof that it is.

\subsection*{Direct Numerical Simulation (DNS)}
DNS solves on a uniform grid resolving the smallest dissipative scales. The state vector has $\mathcal{O}(N_{\text{fine}}^d)$ degrees of freedom, and the explicit CFL condition forces $\Delta t \sim 1/N_{\text{fine}}$.
\begin{itemize}
    \item \textbf{Cost per step:} $T_{\text{DNS}}^{\text{step}} = \mathcal{O}(N_{\text{fine}}^d)$.
    \item \textbf{Cost per unit simulated time:} $T_{\text{DNS}}^{\text{unit}} = \mathcal{O}(N_{\text{fine}}^{d+1})$, obtained by multiplying by the number of steps $\mathcal{O}(N_{\text{fine}})$.
\end{itemize}
Invoking Kolmogorov scaling $N_{\text{fine}} \sim \mathrm{Re}^{3/4}$ in 3D~\cite{pope2000} then gives $T^{\text{step}}_{\text{DNS}} = \mathcal{O}(\mathrm{Re}^{9/4})$ and $T^{\text{unit}}_{\text{DNS}} = \mathcal{O}(\mathrm{Re}^{3})$. The often-quoted $\mathrm{Re}^{9/4}$ figure refers to the cost per step; the cost to simulate a fixed physical time window scales as $\mathrm{Re}^{3}$ in 3D.

\subsection*{Classical Adaptive Mesh Refinement (AMR)}
Let $S_{\text{class}} \in (0,1]$ denote the sparsity factor (fraction of active cells). The classical AMR cost per step splits into a detection stage and a physics stage:
\begin{equation}
    T_{\text{AMR}}^{\text{step}} = \underbrace{\mathcal{O}(N_{\text{fine}}^d)}_{\text{Gradient Scan}} + \underbrace{\mathcal{O}(S_{\text{class}} \cdot N_{\text{fine}}^d)}_{\text{Physics Integration}} + \mathcal{O}(\text{Tree Overhead}).
\end{equation}
Two practical limitations inflate $S_{\text{class}}$: (i) gradient triggers are \emph{reactive} and detect instabilities only after growth, and (ii) they are sensitive to noise, forcing large safety buffers.

\subsection*{Q-HAS Complexity Analysis}
Q-HAS inverts the data dependency: rather than scanning fine gradients bottom-up, it queries coarse topological features top-down through hierarchical VQA calls.

\paragraph{A. Cost of a single VQA call.}
The VQA executes on fixed-size patches ($N_{\text{lim}} \times N_{\text{lim}}$), and its cost is
\begin{equation}
    T_{\text{VQA}} = \mathcal{O}\bigl( K_{\text{opt}}(\epsilon) \cdot N_{\text{shots}}(\epsilon) \cdot p \cdot N_{\text{lim}}^2 \bigr),
\end{equation}
where $K_{\text{opt}}(\epsilon)$ is the number of COBYLA iterations to reach tolerance $\epsilon$, and $N_{\text{shots}}(\epsilon) = \mathcal{O}(\epsilon^{-2})$ is the number of shots per energy estimate to achieve variance $\epsilon^2$. Since neither $K_{\text{opt}}$ nor $N_{\text{shots}}$ has a problem-size-independent upper bound that is tight in practice, we make the working assumption that at fixed target precision $\epsilon$ the per-patch cost $T_{\text{VQA}}(\epsilon)$ is $\mathcal{O}_{N_{\text{fine}}}(1)$, i.e.\ constant in the \emph{global} resolution. This is an empirical heuristic (valid if $K_{\text{max}}$ is capped), not a theoretical guarantee: QAOA lacks an unconditional convergence proof at finite $p$ \cite{farhi2014,zhou2020qaoa}, and the explicit $\epsilon$-dependence of the overall cost is $\mathcal{O}(\epsilon^{-2})$ from shots alone.

\paragraph{B. Detection overhead.}
Even if flux extraction is performed by a full gradient scan $\mathcal{O}(N_{\text{fine}}^d)$, the dominant cost in high-Reynolds MHD is the physics integration (Riemann solvers, magnetic induction), not the error estimation. Performance therefore depends on the sparsity factor $S_{\text{quant}}$ achieved by the quantum-guided refinement. For any quantum advantage to exist, one must have $S_{\text{quant}} < S_{\text{class}}$, which is an empirical question addressed in Chapter~\ref{chap:results}.

\paragraph{C. Tree-search cost.}
The hierarchical BFS visits at most $\mathcal{O}(S_{\text{quant}} \cdot k^j)$ patches at depth $j$, where $k = N_{\text{lim}}^d$ is the branching factor. Summing over depth $j = 0, \dots, D_{\max}$ with $D_{\max} = \lceil \log N_{\text{fine}} / \log N_{\text{lim}} \rceil$ gives a geometric series
\begin{equation}
    \sum_{j=0}^{D_{\max}} k^{j} = \frac{k^{D_{\max}+1} - 1}{k-1} = \mathcal{O}\bigl(k^{D_{\max}}\bigr) = \mathcal{O}(N_{\text{fine}}^d).
\end{equation}
The constant hidden by the $\mathcal{O}(\cdot)$ is $k/(k-1)$, which is close to $1$ for $k \ge 2$. The total tree-search cost is therefore
\begin{equation}
    T_{\text{tree}} = \mathcal{O}\bigl(S_{\text{quant}} \cdot N_{\text{fine}}^d \cdot T_{\text{VQA}}\bigr),
\end{equation}
and the overall Q-HAS cost per step is
\begin{equation}
    T_{\text{Q-HAS}}^{\text{step}} \approx \underbrace{\mathcal{O}\bigl( S_{\text{quant}} \cdot N_{\text{fine}}^d \cdot T_{\text{VQA}} \bigr)}_{\text{Quantum Tree Search}} + \underbrace{\mathcal{O}\bigl( S_{\text{quant}} \cdot N_{\text{fine}}^d \bigr)}_{\text{Sparse Physics Integration}}.
\end{equation}

\subsection*{The Spatio-Temporal Trade-off}

An operational parameter is the interrogation frequency $\nu_{\text{scan}} \in (0,1]$: the fraction of physics steps at which the quantum layer is activated. The amortised per-step cost is
\begin{equation}
    T_{\text{Q-HAS}}^{\text{avg}} \approx \mathcal{O}\bigl(S_{\text{quant}} \cdot N_{\text{fine}}^d \cdot (T_{\text{physics}} + \nu_{\text{scan}} \cdot T_{\text{VQA}})\bigr).
\end{equation}
Comparing to classical AMR with $\nu_{\text{class}} = 1$, a net cost reduction requires
\begin{equation}
    \frac{S_{\text{quant}}}{S_{\text{class}}} < \frac{T_{\text{physics}}}{T_{\text{physics}} + \nu_{\text{scan}} \cdot T_{\text{VQA}}}.
    \label{eq:advantage-condition}
\end{equation}
This inequality identifies two levers --- spatial sparsity ($S_{\text{quant}} \ll S_{\text{class}}$) and temporal sparsity ($\nu_{\text{scan}} \ll 1$) --- but the empirical results of Chapter~\ref{chap:results} do not currently satisfy \eqref{eq:advantage-condition}: at $N = 256$, $\mathrm{Re} = 400$, $S_{\text{quant}}/S_{\text{class}}$ is close to or above $1$ on three of the four main scenarios, and $T_{\text{VQA}}$ on simulator is several orders of magnitude larger than $T_{\text{physics}}$. The asymptotic regime in which the condition holds is therefore not reached in this work; scaling studies at higher $\mathrm{Re}$ and with larger registers $N_{\text{lim}}$ are needed to test it.

\subsection*{Future Scalability: The Boundary Sampling Regime}
With small quantum registers, coarse fluxes still require $\mathcal{O}(N_{\text{fine}}^d)$ volume averaging. Larger registers would enable a scaling shift via the divergence theorem: internal anomalies manifest as flux imbalances at coarse block boundaries, and the detection complexity could in principle reduce from $\mathcal{O}(N_{\text{fine}}^d)$ to $\mathcal{O}(N_{\text{fine}}^{d-1})$. This regime is not reached here and is discussed as a future direction only.


\chapter{Training Hyperparameters}
\label{chap:training}

Q-HAS relies on phenomenological hyperparameters bridging classical plasma states and quantum operators. These parameters depend on the plasma regime ($Re$, $Rm$). Optimisation uses a \textit{Simulation-in-the-Loop} strategy with the Optuna Bayesian framework~\cite{akiba2019}.

\section{The Multi-Objective Cost Function}
The optimisation seeks to minimise a hybrid loss function $\mathcal{J}$ that penalises both physical error and computational resource usage:
\begin{equation}
    \mathcal{J}(\Theta) = \frac{\mathcal{E}_{\text{phys}}(\Theta) + \lambda \cdot \mathcal{S}(\Theta)}{1 + \lambda}
    \label{eq:combined-score}
\end{equation}
where $\Theta$ represents the set of trainable hyperparameters:\\ $\{\beta, \text{threshold}_{\text{amr}}, \gamma_{\text{hydro}}, \gamma_{\text{mag}}, \beta_{\text{grad}}, \beta_{\text{curl}}, \beta_{\text{xpoint}}, \kappa, w_{z,\text{frac}}\}$.

\begin{itemize}
    \item \textbf{Asymmetric Physical Fidelity ($\mathcal{E}_{\text{phys}}$):} We quantify accuracy using a weighted Relative $L_2$ Norm across five critical MHD fields: fluid velocity ($v_x, v_y$), magnetic field ($B_x, B_y$), and current density ($J_z$). We apply a spatially-varying weight map that combines both current density and vorticity:
    \begin{equation}
        w(x,y) = 1 + 0.5 \cdot \frac{|J_z(x,y)| / \langle|J_z|\rangle + |\omega_z(x,y)| / \langle|\omega_z|\rangle}{2}
    \end{equation}
    This ensures that missing an instability (false negative) in high-current or high-vorticity zones is penalised more heavily than over-refining a quiet region. Using both indicators ensures correct penalisation even in scenarios where only one is active (e.g., Lamb--Oseen vortex has high $|\omega_z|$ but low $|J_z|$).

    \item \textbf{Computational Sparsity ($\mathcal{S}$):} This represents the average fraction of the grid refined by the algorithm across all hybrid steps: $\mathcal{S} = \text{avg\_pixels\_refined} / N^2$.

    \item \textbf{The Trade-off Parameter ($\lambda$):} Set to $\lambda = 0.4$ in the Phase~1 training runs used in this report. This value imposes a meaningful sparsity constraint while keeping physical accuracy as the primary optimisation target; values closer to $\lambda = 0$ push the optimiser towards full refinement, while values closer to $\lambda = 1$ collapse the solution onto the coarse grid.
\end{itemize}

\section{Three-Phase Progressive Training}
\label{sec:multi-scenario-training}
The optimisation was designed as a progressive three-phase schedule. In the work reported here, only Phase~1 was executed with split sensitivities (5 free quantum parameters) (not enough time for the analysis of others); the threshold was frozen from classical training. A DNS reference trajectory is precomputed once per scenario and reused across all Optuna trials, halving per-trial compute cost by skipping the witness solver.

\begin{enumerate}
    \item \textbf{Phase 1 (executed). Gate calibration with split sensitivities.} Trains 5 quantum parameters ($\beta$, $\sigma$, $\beta_{\text{curl}}$, $\beta_{\text{xpoint}}$, $w_{z,\text{frac}}$) on four scenarios simultaneously, each targeting a specific anomaly type:
    \begin{itemize}
        \item \textbf{Kelvin--Helmholtz:} Pure shear layers without current sheets, which calibrates $g_{\text{strain}}$ and the gradient coupling.
        \item \textbf{Harris Tearing:} Current sheet with reconnection sites, which calibrates $g_{\text{mag}}$ and selective detection.
        \item \textbf{Orszag--Tang:} Mixed turbulence with all anomaly types, providing a composite benchmark.
        \item \textbf{MHD Rotor:} Rotating field configuration generating vorticity from $t = 0$.
    \end{itemize}
    The frozen parameters ($\text{threshold}_{\text{amr}} = 0.1496$ from classical training, $\gamma_{\text{hydro}} = 2.0$, $\gamma_{\text{mag}} = 0.5$, $\kappa = 10.0$) were calibrated in earlier training runs or set by physical reasoning. The composite loss is the average over the four scenarios:
    \begin{equation}
        \mathcal{J}_{\text{composite}} = \frac{1}{4}\left(\mathcal{J}_{\text{KH}} + \mathcal{J}_{\text{Tearing}} + \mathcal{J}_{\text{OT}} + \mathcal{J}_{\text{Rotor}}\right)
    \end{equation}

    \item \textbf{Phase 2 (defined, not executed). Multi-scenario generalisation.} Would train all parameters on an expanded scenario set including Lamb--Oseen vortex and island coalescence.

    \item \textbf{Phase 3 (defined, not executed). Validation and fine-tuning on Orszag--Tang.} Would validate generalisation to the combined Orszag--Tang flow at $Re = Rm = 800$.
\end{enumerate}

\subsection{Hyperparameter Reference}

\begin{table}[H]
\centering
\caption{Trainable hyperparameters and their search ranges. Five quantum parameters are trained; four are frozen from prior calibration.}
\label{tab:hyperparams}
\begin{tabular}{llccc}
\toprule
\textbf{Parameter} & \textbf{Description} & \textbf{Range} & \textbf{Best} & \textbf{fANOVA} \\
\midrule
$\beta_{\text{curl}}$    & $T_{rcf}$ sensitivity for plaquette $ZZZZ$ terms & $[0.1, 10.0]$  & 4.27 & 44\% \\
$\beta_{\text{xpoint}}$  & $T_{rcf}$ sensitivity for X-point $ZZZZ$ terms   & $[0.1, 10.0]$  & 2.39 & 21\% \\
$\sigma$                 & Uncertainty width for ZZ coupling Gaussian        & $[0.01, 0.30]$ & 0.023 & 18\% \\
$\beta$                  & Phase encoding gain ($\phi$ sensitivity)           & $[0.5, 10.0]$  & 9.94 & 11\% \\
$w_{z,\text{frac}}$      & Z-bias fraction ($\alpha_z$ scaling)               & $[1.0, 200.0]$ & 10.40 & 6\% \\
\midrule
\multicolumn{5}{l}{\textit{Frozen parameters}} \\
$\text{threshold}_{\text{amr}}$ & AMR recursion threshold (from classical training) & -- & 0.1496 & -- \\
$\gamma_{\text{hydro}}$  & $f$-gate log-growth for $Re$                      & -- & 2.0 & -- \\
$\gamma_{\text{mag}}$    & $f$-gate log-growth for $Rm$                      & -- & 0.5 & -- \\
$\kappa$                 & Leaky sigmoid steepness                            & -- & 10.0 & -- \\
\bottomrule
\end{tabular}
\end{table}

\textbf{Key insight from fANOVA ("percentage of importance of each parameter to minimize the loss function" :} $\beta_{\text{curl}}$ (44\%) dominates, meaning the QAOA's main lever is magnetic reconnection detection via the curl-based $ZZZZ$ terms. $w_{z,\text{frac}}$ at only 6\% confirms that $ZZ$ spatial correlations barely influence the training score at $2\times 2$ VQA resolution. The trained $\sigma = 0.023$ is very tight, concentrating ZZ coupling on a narrow band around the threshold.

% ==========================================
% CHAPTER: COMPUTATIONAL IMPLEMENTATION
% ==========================================
\chapter{Computational Implementation}
\label{chap:implementation}

The complete source code is publicly available at \url{https://github.com/armandld/BA_Proj}.

\section{The Hybrid Software Stack}
The framework is organised into four interacting layers, all implemented in Python 3.11:
\begin{itemize}
    \item \textbf{Classical MHD Solver} (\texttt{src/Simulation/}): The core physics engine, implemented using \texttt{NumPy}~\cite{numpy2020} and \texttt{SciPy}~\cite{scipy2020} for high-performance array operations. The solver (\texttt{solver.py}) handles spatial discretisation via FD4, temporal integration via RK4 with CFL-adaptive time stepping, and divergence-free projection via spectral FFT. Grid infrastructure (\texttt{grid.py}) provides a $[0, 2\pi]^2$ periodic domain with gradient, divergence, Laplacian, and FFT-based Poisson operators. The adaptive mesh refinement engine (\texttt{refinement.py}) implements both the quantum-guided (level-by-level BFS) and classical-baseline scanning. Two stepping modes are provided: \texttt{step\_full()} for full DNS evolution and \texttt{step\_layered()} for hierarchical AMR.
    \item \textbf{Physics-to-Quantum Mapping} (\texttt{src/Simulation/PhysToAngle.py}, \texttt{HamiltParams.py}): The \texttt{AngleMapper} serves a dual purpose. First, it computes the \textbf{classical multi-indicator score} (Equation~\ref{eq:classical-score}) combining vorticity, divergence, $|J_z|$, and the L\"ohner estimator on $|\mathbf{B}|$, which determines the polar angle $\theta$ (Equation~\ref{eq:theta-encoding}). Second, it computes the \textbf{stress flux} $\Phi$ using a shock-diode logic (separating compression from shear) and maps its temporal derivative to the phase angle $\phi$ (Equation~\ref{eq:phi-encoding}). The stress flux is also used by the Hamiltonian. The \texttt{PhysicalMapper} computes all Hamiltonian coefficient arrays using the decoupled $f \times g \times T_{rcf}$ architecture (Section~\ref{sec:hamilt-architecture}), with absolute physical non-dimensionalisation (fixed critical thresholds $Re_{\text{crit}} = 1$, $Rm_{\text{crit}} = 1$, $Q_{\text{crit}} = 2$, $J_{\text{crit}} = 1$).
    \item \textbf{Quantum Routine} (\texttt{src/VQA/}): The VQA pipeline chains five stages, each implemented as a separate module:
    \begin{enumerate}
        \item \textit{Mapping} (\texttt{mapping.py}): Constructs the full QAOA circuit by composing the warm-start initial state (Section~\ref{sec:gate-level}) with the Qiskit \texttt{QAOAAnsatz} built from the cost Hamiltonian. For periodic patches (depth~0), \texttt{create\_period\_hamiltonian} is used; for bounded patches (depth~$>0$), \texttt{create\_bounded\_hamiltonian} extracts core-vs-halo angles and replaces boundary qubits with classical $\langle Z \rangle = \cos(\theta)$ expectation values.
        \item \textit{Optimisation} (\texttt{optimize.py}): Transpiles the parametrised circuit for the target backend using Qiskit's preset pass manager. Statevector backends use \texttt{opt\_level=0} (no routing required), while noisy backends use higher optimisation levels.
        \item \textit{Execution} (\texttt{execute.py}): Runs the COBYLA-based minimisation loop. The initial parameters follow the linear-ramp strategy (Section~\ref{sec:initialization}): all mixer angles $\vec{\Omega}$ are initialised to zero and the structural angles $\vec{\Gamma}$ follow the ramp $\Gamma_k = (2k / p(p+1)) \cdot \pi / E_{\max}$, where $E_{\max}$ is the sum of absolute Hamiltonian weights. Energy evaluations use \texttt{EstimatorV2}, and the final optimised circuit is measured with \texttt{SamplerV2}. Alternative optimisers (\texttt{L-BFGS-B}, \texttt{Powell}) are also supported, with method-specific bound and tolerance settings.
        \item \textit{Post-processing} (\texttt{postprocess.py}): Converts the sampled bitstring distribution into per-qubit marginal probabilities $P(q_k = 1)$ by iterating over all measured bitstrings (using Qiskit's little-endian bit-ordering convention, reversed to natural qubit order) and accumulating weighted counts.
        \item \textit{Runtime} (\texttt{runtime.py}): A \texttt{VQARuntime} singleton manages shared\\ \texttt{EstimatorV2}/\texttt{SamplerV2} primitives and an ansatz cache keyed by \texttt{(num\_qubits, period\_bound, reps)}. This avoids repeated primitive instantiation and expensive ansatz decomposition across the many VQA calls in a single simulation.
    \end{enumerate}
    Quantum circuits are constructed with \texttt{Qiskit}~\cite{qiskit2024} and executed on the \texttt{Qiskit-Aer} statevector simulator. The code also supports execution on IBM's \texttt{FakeFez} noisy simulator and, when available, on real IBM quantum hardware via the \texttt{--mode hardware} command-line flag with dynamical decoupling and gate twirling for error mitigation.
    \item \textbf{Training Orchestration} (\texttt{src/TrainHyperParam\_v2.py}): \texttt{Optuna}~\cite{akiba2019} manages the 3-phase Bayesian optimisation (Section~\ref{sec:multi-scenario-training}) with a \texttt{MedianPruner} (15 startup trials, minimum 5 completed before pruning) to terminate unpromising trials early. Phase~2 uses a composite multi-scenario approach where the loss is averaged over isolated scenarios (KH, vortex, tearing, coalescence), ensuring each $f \times g \times T_{rcf}$ gate parameter is calibrated on the anomaly type it is designed to detect. Distributed training is supported via three storage backends: PostgreSQL (for Google Colab with Neon), JournalFileStorage (for HPC clusters with shared NFS), and local SQLite (default).
    \item \textbf{Validation Suite} (\texttt{tests/}): A multi-tier test suite implemented using \texttt{pytest} validates solver convergence, AMR consistency, and VQA anomaly detection across eight study cases.
\end{itemize}

\subsection{The AMR Solver}

The AMR solver (\texttt{step\_layered}) implements a three-phase defect-correction architecture:
\begin{enumerate}
    \item \textbf{Phase 1 ,  Global coarse step:} The entire domain is downsampled by the coarse factor ($\text{target\_dim}^{D_{\max}}$), evolved by one RK4 step at coarse resolution, and the delta is upsampled back to full resolution (using bicubic periodic interpolation). This provides a cheap baseline correction for the whole domain.
    \item \textbf{Phase 2 ,  Per-patch finer resolution:} For each active patch, the local region is extracted with padding (8~pixels for the FD4$\times$RK4 stencil), downsampled to the patch's local resolution (determined by its depth), and evolved by a local RK4 step. Deeper patches (higher depth) have smaller \texttt{local\_factor}, yielding higher resolution.
    \item \textbf{Phase 3 ,  Global stabilisation:} The FFT-based divergence-free projection is applied to the full domain (identical to \texttt{step\_full}), ensuring incompressibility.
\end{enumerate}
\textbf{Convergence guarantee:} When all patches are at $D_{\max}$ (where $\texttt{local\_factor} = 1$), every patch's local RK4 step runs at full DNS resolution. The result is mathematically identical to \texttt{step\_full()}, up to floating-point precision ($< 10^{-10}$), as verified by the test suite.

\section{Numerical Methods: FD4 and RK4}
Numerical schemes are chosen for AMR compatibility:
\begin{itemize}
    \item \textbf{Spatial Discretisation:} 4th-order centred \textbf{Finite Differences (FD4)}. The local stencil allows independent solving on arbitrary rectangular sub-patches, essential for patch-based AMR (unlike global spectral methods).

    The Finite Volume Method (FVM) would be preferable: it derives from integral conservation laws, maintaining conservation exactly across resolution changes. FD requires interpolation at AMR boundaries, introducing non-conservative errors. FVM also acts as a natural low-pass filter via surface integrals ($\oint_{\partial\Omega} \mathbf{F} \cdot d\mathbf{S}$), whereas FD is sensitive to single-pixel noise. A full FVM implementation was beyond scope (see Chapter~\ref{chap:limitations}).

    
    \item \textbf{Temporal Integration:} 4th-order \textbf{Runge--Kutta (RK4)}. Centred FD schemes are non-dissipative, so RK4's stability region (function $R_4(z)= \sum_{k=0}^4 z^k / k!$, stable for $|R_4(z)|<1$) must cover the imaginary axis (approximately $[-2.82\,i,\; 2.82\,i]$).

    This is essential because FD4 applied to advection produces purely imaginary eigenvalues, requiring a time integrator with imaginary-axis stability coverage that RK2 lacks.

    \item \textbf{Incompressibility:} Two divergence-free constraints must hold at every instant:
    \begin{itemize}
        \item $\nabla \cdot \mathbf{v} = 0$ (mass conservation, constant density).
        \item $\nabla \cdot \mathbf{B} = 0$ (Gauss’s law, no magnetic monopoles).
    \end{itemize}
    RK4 and FD4 do not preserve these automatically; truncation errors create non-physical divergence. The \textbf{Projection Method} (Chorin~\cite{chorin1968}) enforces these constraints: after each step, an FFT-based spectral Poisson solver projects $\mathbf{v}$ and $\mathbf{B}$ onto divergence-free subspaces.

    For velocity: $\nabla^2 \varphi = \nabla \cdot \vec{v}^*$ (where $\vec{v}^*$ is the uncorrected field), then $\vec{v}_{\text{new}} = \vec{v}^* - \nabla \varphi$. The same applies to $\mathbf{B}$. FFT-based solving has $\mathcal{O}(N\log N)$ complexity.
\end{itemize}

\section{Software Validation and Test Suite}
\label{sec:test-suite}

A multi-tier test suite validates solver integrity and VQA anomaly detection sensitivity.

\subsection{Numerical Convergence and Solver Integrity}

Five convergence tests validate the AMR engine: exact convergence (\texttt{step\_layered} matches \texttt{step\_full} to $< 10^{-10}$ at max depth), normalisation of the patch cost metric, physics hierarchy (L2 error monotonically decreasing with depth), the \texttt{compute\_local\_factor} helper, and the corrected \texttt{max\_depth} subdivision formula.

\subsection{VQA Anomaly Study Cases}
\label{sec:vqa-study-cases}

Eight study cases validate each Hamiltonian term against synthetic plasma states: shear detection ($ZZ$), circulation ($ZZZZ$), X-point reconnection ($ZZZZ$), phase encoding ($\phi$), coefficient sign verification (all positive for correct minimisation), energy minimisation proof ($E_{\text{QAOA}} < E_{\text{all-}|0\rangle}$), cross-anomaly isolation (ensuring sensor independence), and combined anomalies (verifying simultaneous activation).

\section{Physical Scenarios}
\label{sec:physical-scenarios}

The solver supports eight MHD instability scenarios: Orszag--Tang vortex (mixed turbulence benchmark), Kelvin--Helmholtz (shear-driven rolling vortices), Harris tearing (current sheet with reconnection), magnetic twist (``silent'' instability invisible to classical AMR), noisy uniform field (noise immunity test), double tearing (frustrated current sheet interaction), Lamb--Oseen vortex (isolated vorticity), island coalescence (merging magnetic islands with X-points), and MHD rotor (rotating field generating vorticity from $t = 0$). The first four serve as the primary training scenarios; the remainder are used for isolated-anomaly validation.

\chapter{Computational Results}
\label{chap:results}

This chapter presents the computational benchmarks that probe whether Q-HAS offers measurable advantage over a classically-trained gradient AMR baseline. The evaluation is deliberately conservative: both pipelines share the identical classical MHD solver, the same hierarchical BFS refinement engine, the same TTL mechanism, and the same domain, resolution and random seeds. The only difference is the decision variable: the classical pipeline thresholds the four-indicator score $s$ of Equation~\ref{eq:classical-score}, while the quantum pipeline thresholds the QAOA-refined probability map $P(|q\rangle = |1\rangle)=|\langle q|1\rangle|^2$. Each method uses its \textit{own} optimally-trained decision threshold, calibrated independently during training. Physical parameters are $Re = 400$, $Rm = 400$, $\eta = 0.001$, $c_s = 1.0$, $B_{z,\text{guide}} = 0.1$.

All figures in this chapter were generated by the \texttt{figures\_code/} scripts in the public repository and saved under \texttt{figures/}. The numerical values quoted in the text are taken directly from the figure log files.

\section{Training Results}
\label{sec:training-results}

\begin{table}[H]
\centering
\caption{Training outcome. Q-HAS achieves a marginal 0.66\% advantage in composite loss.}
\label{tab:training-outcome}
\begin{tabular}{lccc}
\toprule
\textbf{Method} & \textbf{Best Score} & \textbf{Best Trial} & \textbf{Trials Run} \\
\midrule
Q-HAS     & \textbf{0.2134} & \#4  & 170+ \\
Classical & 0.2148          & \#42 & 125 \\
\midrule
\textbf{Advantage} & \textbf{0.66\%} & -- & -- \\
\bottomrule
\end{tabular}
\end{table}

The quantum advantage is real but marginal. Optuna explored the parameter space widely ($w_{z,\text{frac}}$ from 7 to 168) without strong convergence, suggesting the classical threshold is already near-optimal at this scale ($N = 256$, $Re = Rm = 400$, $2\times 2$ VQA grid).

\subsection{Pareto Frontiers}
\label{sec:pareto}

Figure~\ref{fig:pareto-all} overlays the Pareto fronts of the quantum and classical trials in the $(\mathcal{E}_{\text{phys}}, \mathcal{S})$ plane for all lambda values $\lambda = 0.4$ (a balanced value to avoid over refinement without letting a very coarse grid doing a sufficient, approximated simulation).

\begin{figure}[H]
    \centering
    \includegraphics[width=0.92\textwidth]{figures/fig0_pareto_all_combined.png}
    \caption[Pareto frontier: all scenarios combined]{Combined Pareto frontier across the four Phase~1 scenarios. The operating point at $\lambda = 0.40$ sits in the knee of the Pareto curve where further cost reduction yields diminishing physics fidelity.}
    \label{fig:pareto-all}
\end{figure}

% ===========================
\section{Per-Scenario Comparison}
\label{sec:per-scenario-results}

\subsection{Comprehensive Comparison}

\begin{figure}[H]
    \centering
    \includegraphics[width=\textwidth]{figures/fig4_comprehensive_comparison.png}
    \caption{Per-scenario captured fraction, precision, recall, and compute ratio for Q-HAS vs.\ classical AMR.}
    \label{fig:comprehensive}
\end{figure}

Table~\ref{tab:per-scenario} summarises the numerical results.

\begin{table}[H]
\centering
\caption{Per-scenario comparison: Q-HAS vs Classical AMR ($N = 256$, 2$\times$2 VQA, 8 qubits).}
\label{tab:per-scenario}
\begin{tabular}{lcccccccc}
\toprule
 & \multicolumn{2}{c}{\textbf{Captured}} & \multicolumn{2}{c}{\textbf{Precision}} & \multicolumn{2}{c}{\textbf{Recall}} & \multicolumn{2}{c}{\textbf{Compute}} \\
\textbf{Scenario} & Q-HAS & Class. & Q-HAS & Class. & Q-HAS & Class. & Q-HAS & Class. \\
\midrule
KH      & 0.955 & \textbf{0.968} & \textbf{0.705} & 0.683 & 0.975 & \textbf{1.000} & \textbf{0.754} & 0.846 \\
Tearing & 0.953 & \textbf{0.984} & 1.000 & 1.000 & 0.860 & \textbf{0.865} & 0.625 & 0.625 \\
Rotor   & \textbf{0.976} & 0.971 & 0.983 & \textbf{1.000} & 0.661 & \textbf{0.834} & 0.398 & \textbf{0.273} \\
OT      & 0.695 & \textbf{0.729} & 0.700 & \textbf{0.777} & 0.277 & \textbf{0.441} & 0.623 & \textbf{0.621} \\
\bottomrule
\end{tabular}
\end{table}

\textbf{Kelvin--Helmholtz (best case for Q-HAS):} Q-HAS uses 10.9\% less compute (0.754 vs 0.846) with only marginal loss in captured fraction ($-1.3\%$) and better precision ($+3.2\%$). The KH instability has a simple, elongated structure; Q-HAS avoids over-refinement in the bulk.

\textbf{Harris Tearing:} Both methods achieve high captured fraction ($>0.95$) and precision (1.000). Compute ratio is identical. The simple sheet-like structure of the current sheet is easy for both methods.

\textbf{MHD Rotor (worst case for Q-HAS):} Q-HAS has a significant recall deficit ($0.661$ vs $0.834 = -20.7\%$) and uses more compute ($0.398$ vs $0.273 = +45.8\%$). The rotor's circular geometry with strong gradients at the interface leads the QAOA corrections to flag a wider buffer than the classical detector, inflating $S_{\text{quant}}$ while missing part of the core.

\textbf{Orszag--Tang (complex turbulence):} Classical wins on every metric. Q-HAS has lower captured fraction ($-4.7\%$), lower precision ($-9.9\%$), and much lower recall ($0.277$ vs $0.441 = -37.2\%$). The distributed, multi-scale turbulent structure is the hardest for both methods.

\textbf{Overall assessment:} Q-HAS shows a genuine efficiency advantage \textit{only} on Kelvin--Helmholtz. On Rotor and Orszag--Tang, Q-HAS is clearly worse. On Tearing, it is neutral. There is no consistent quantum advantage across scenarios.

\subsection{Statistical Validation}

\begin{figure}[H]
    \centering
    \includegraphics[width=\textwidth]{figures/fig6_statistical_validation.png}
    \caption{Bootstrap confidence intervals and permutation tests for the captured-fraction advantage $\Delta = F_{\text{Q-HAS}} - F_{\text{classical}}$, computed over 10 independent seeds.}
    \label{fig:statistical}
\end{figure}

\begin{table}[H]
\centering
\caption{Statistical significance of the captured-fraction advantage $\Delta = F_{\text{Q-HAS}} - F_{\text{classical}}$ (one-sided paired permutation test + bootstrap, 10 seeds). Cohen's $d$ is computed on $\Delta$: a positive value indicates Q-HAS is better on that metric, a negative value indicates classical is better. Note that with $n=10$ seeds the effect-size estimates have wide sampling error and the listed $p$-values are \emph{not} corrected for the 4 scenarios tested (Bonferroni-corrected threshold for $\alpha = 0.05$ would be $p \le 0.0125$).}
\label{tab:statistical}
\begin{tabular}{lcccc}
\toprule
\textbf{Scenario} & \textbf{Direction (favoured)} & \textbf{$p$-value (one-sided)} & \textbf{Cohen's $d$ on $\Delta$} \\
\midrule
KH      & Q-HAS      & 0.008 & $+0.90$ \\
Tearing & Classical  & 0.001 & $-3.41$ \\
Rotor   & Classical (n.s.) & 0.974 & $-0.64$ \\
OT      & Classical (n.s.) & 0.999 & $-1.01$ \\
\bottomrule
\end{tabular}
\end{table}

\paragraph{Interpretation and caveats.}
The script used to generate Fig.~\ref{fig:statistical}(A) uses a one-sided permutation test that tests only the alternative ``Q-HAS better'', and labels the opposite direction as ``n.s.''; we have relabelled the Direction column so that the sign of Cohen's $d$ and the favoured method agree. After this relabelling, only KH survives as a genuine positive for Q-HAS ($p = 0.008$, $d = +0.90$). Tearing shows a highly significant regression in the Q-HAS direction ($p = 0.001$, $d = -3.41$), which the one-sided protocol cannot detect as significant: a two-sided test would flag it. Rotor and OT are consistent with classical being better but do not reach significance at $n=10$. None of these $p$-values survive Bonferroni correction for 4 scenarios except KH (marginally) and Tearing. Three further caveats apply: (i) $n=10$ seeds is underpowered for $|d|>1$ effects --- the very large $|d|=3.41$ on Tearing should be reproduced at $n \ge 30$ before being taken at face value; (ii) the results compare to a classical baseline implemented in this code base, not to a production AMR code such as FLASH or PLUTO; (iii) the test is applied to captured fraction only, not to the full joint (accuracy, sparsity) objective.

% ===========================
\section{QAOA Mechanism Analysis}

\subsection{Detailed QAOA Corrections}

\begin{figure}[H]
    \centering
    \includegraphics[width=\textwidth]{figures/fig5_qaoa_detailed_analysis.png}
    \caption{Per-cell QAOA correction $\delta = P_{\text{QAOA}}(|1\rangle) - s_{\text{classical}}$, Hamiltonian energy decomposition, and hierarchical summary.}
    \label{fig:qaoa-detail}
\end{figure}

Corrections are \textbf{non-zero} everywhere, proving the QAOA circuit is functional. The corrections are structured (vary by scenario, depth, cell position) and are not random noise. However, $H(Z)$ coefficients dominate after downsampling to the $2\times 2$ VQA grid; $C(ZZ)$ and $K(ZZZZ)$ terms are visible but small relative to $H(Z)$. The QAOA is therefore Z-dominated: the spatial correlations are a perturbative correction, not the primary signal.

\subsection{Hamiltonian Design Visualisation}

\begin{figure}[H]
    \centering
    \includegraphics[width=\textwidth]{figures/fig11_hamiltonian_design.png}
    \caption{Hamiltonian design: classical score, uncertainty-weighted ZZ coupling (always-on vs.\ trained $\sigma = 0.023$), uncertainty weight field, and Z-bias for KH (top) and OT (bottom).}
    \label{fig:hamiltonian-design}
\end{figure}

For KH, the ZZ coupling is effectively \textbf{zero} ($\sim 10^{-6}$): the Gaussian gate $\exp(-((\bar{s} - 0.15)/0.023)^2)$ kills all coupling because KH scores are far from the threshold. Only for OT does the ZZ coupling activate meaningfully along vortex filaments. This explains why $\sigma = 0.023$ limits quantum advantage: for scenarios with bimodal score distributions, the Hamiltonian is effectively null.

\subsection{Sigma Ablation Study}

\begin{figure}[H]
    \centering
    \includegraphics[width=\textwidth]{figures/fig13_sigma_ablation.png}
    \caption{Ablation study: ZZ/Z ratio, active edge fraction, coefficient magnitudes, and score distance from threshold as functions of $\sigma$.}
    \label{fig:sigma-ablation}
\end{figure}

At the trained $\sigma = 0.023$: only 5--20\% of edges have active ZZ coupling (uncertainty weight $> 0.1$), and for KH the ZZ magnitude is $\sim 10^{-7}$ while Z stays at $\sim 10^{-3}$, a factor of $10^4$ weaker. The ablation proves that $\sigma = 0.023$ is too narrow for the current score distributions. A broader $\sigma$ (0.10--0.15) would activate more edges, but Optuna chose $\sigma = 0.023$ because it minimised the training loss, suggesting that activating more edges does not help (the ZZ corrections hurt as often as they help).

% ===========================
\section{Decision-Level Diagnostics}

\subsection{Decision Flip Analysis}
\label{sec:decision-flip}

\begin{figure}[H]
    \centering
    \includegraphics[width=\textwidth]{figures/fig15_decision_flip_analysis.png}
    \caption{THE diagnostic figure: classical score vs QAOA probability for every BFS cell, threshold proximity, flips by depth, correction--threshold comparison, accuracy summary, and physical mechanism (WHY) analysis.}
    \label{fig:flip-analysis}
\end{figure}

\begin{table}[H]
\centering
\caption{Decision flip analysis ($N = 256$, depth = 5). Classical is more accurate on aggregate.}
\label{tab:flips}
\begin{tabular}{lccccc}
\toprule
\textbf{Scenario} & \textbf{Cells} & \textbf{Flips} & \textbf{Flip Rate} & \textbf{QAOA Correct} & \textbf{Classical Correct} \\
\midrule
KH      & 904 & 125 & 13.8\% & 61.3\% & 65.6\% \\
Tearing & 584 & 109 & 18.7\% & 77.9\% & 81.2\% \\
Rotor   & 240 & 76  & 31.7\% & 67.9\% & 72.1\% \\
OT      & 408 & 145 & 35.5\% & 68.6\% & 70.3\% \\
\midrule
\textbf{Total} & 2136 & 455 & \textbf{21.3\%} & 43.3\% & 46.6\% \\
\bottomrule
\end{tabular}
\end{table}

The flip rate (21.3\%) is higher than expected, showing QAOA corrections \textit{do} change decisions at a meaningful fraction of cells. However, classical is more accurate by 71 cells overall. ZZ energy is 783,000$\times$ higher at flip sites (median 0.0078 vs $\sim 0$ for non-flipped cells), confirming that the ZZ coupling drives corrections at the decision boundary. 44.5\% of corrections have $|\text{correction}| > |\text{distance to threshold}|$ (ratio $> 1.0$); the median ratio is 0.73.

\subsection{Decision Landscape}

\begin{figure}[H]
    \centering
    \includegraphics[width=\textwidth]{figures/fig16_decision_landscape.png}
    \caption{Per-cell scatter plot of classical score vs QAOA probability, coloured by ground-truth correctness.}
    \label{fig:decision-landscape}
\end{figure}

\begin{table}[H]
\centering
\caption{Decision landscape: flip counts and net winner per scenario.}
\label{tab:landscape}
\begin{tabular}{lcccccc}
\toprule
\textbf{Scenario} & \textbf{Cells} & \textbf{Flips} & \textbf{Flip Rate} & \textbf{QAOA Correct} & \textbf{Class.\ Correct} & \textbf{Net Winner} \\
\midrule
KH      & 884 & 137 & 15.5\% & 61.4\% & 64.9\% & Classical by 31 \\
Tearing & 588 & 111 & 18.9\% & 77.7\% & 81.0\% & Classical by 19 \\
Rotor   & 252 & 80  & 31.7\% & 70.2\% & 75.0\% & Classical by 12 \\
OT      & 432 & 144 & 33.3\% & 69.9\% & 67.6\% & \textbf{Q-HAS by 10} \\
\bottomrule
\end{tabular}
\end{table}

OT is the \textit{only} scenario where Q-HAS wins on decision accuracy, consistent with OT having the most active ZZ/ZZZZ coupling and the richest topological structure.

% ===========================
\section{Topological Attribution}
\label{sec:topological-attribution}

\begin{figure}[H]
    \centering
    \includegraphics[width=\textwidth]{figures/fig17_topological_attribution.png}
    \caption{Topological correction attribution: accuracy by ZZ quartile, flip accuracy by dominant Hamiltonian term, and Q-HAS advantage in topology-rich vs smooth regions.}
    \label{fig:topological}
\end{figure}

\begin{table}[H]
\centering
\caption{Flip accuracy by dominant Hamiltonian term. ZZ-driven flips are 76.5\% accurate for Tearing; ZZZZ-driven flips are 66.3\% accurate for OT.}
\label{tab:flip-by-term}
\begin{tabular}{lcccccc}
\toprule
\textbf{Scenario} & \multicolumn{2}{c}{\textbf{ZZ (spatial)}} & \multicolumn{2}{c}{\textbf{ZZZZ (topology)}} & \multicolumn{2}{c}{\textbf{Z (bias)}} \\
 & Acc. & Flips & Acc. & Flips & Acc. & Flips \\
\midrule
KH      & 0.333 & 45 & 0.500 & 14 & 0.431 & 58 \\
Tearing & \textbf{0.765} & 68 & 0.000 & 8  & 0.000 & 41 \\
Rotor   & 0.429 & 14 & 0.491 & 53 & 0.333 & 15 \\
OT      & 0.378 & 45 & \textbf{0.663} & 92 & 0.200 & 5  \\
\bottomrule
\end{tabular}
\end{table}

\begin{table}[H]
\centering
\caption{Q-HAS advantage by topological regime. Topology-rich regions show positive advantage for Tearing ($+3.8$pp) and OT ($+5.3$pp).}
\label{tab:topo-regime}
\begin{tabular}{lcccc}
\toprule
\textbf{Scenario} & \textbf{Topology-rich $\Delta$} & \textbf{Smooth $\Delta$} & \textbf{Q-HAS rich} & \textbf{Classical rich} \\
\midrule
KH      & $-0.042$ & $-0.012$ & 0.322 & 0.363 \\
Tearing & $\mathbf{+0.038}$ & $-0.082$ & 0.729 & 0.692 \\
Rotor   & $-0.057$ & $-0.008$ & 0.631 & 0.689 \\
OT      & $\mathbf{+0.053}$ & $+0.024$ & 0.675 & 0.621 \\
\bottomrule
\end{tabular}
\end{table}

\textbf{Central finding:} In topology-rich regions, Q-HAS shows a positive advantage for Tearing ($+3.8$ percentage points) and OT ($+5.3$pp). These are the two scenarios with the strongest topological features (current-sheet reconnection for Tearing, vortex interactions for OT). In smooth regions, the advantage disappears or reverses.

\textit{Note} (Harris Tearing, Fig.~\ref{fig:topological}): The topological correction analysis reveals that in the Harris Tearing scenario, Q-HAS performs 97 \textit{incorrect} flips against only 50 \textit{correct} flips relative to ground truth. The system tends to add unnecessary refinement or to remove critical areas, suggesting that the X-point detection term ($ZZZZ$) is poorly calibrated at this resolution.

% ===========================
\section{Physical Fidelity and Structural Analysis}

\subsection{Physical Fidelity}

\begin{figure}[H]
    \centering
    \includegraphics[width=\textwidth]{figures/fig7_physical_fidelity.png}
    \caption{L2 error vs DNS, kinetic energy conservation, and enstrophy conservation over 15 AMR-guided evolution steps.}
    \label{fig:physical-fidelity}
\end{figure}

\begin{table}[H]
\centering
\caption{Final L2 error vs DNS after 15 AMR steps ($N = 256$, depth = 5). Classical wins on all scenarios.}
\label{tab:fidelity}
\begin{tabular}{lccc}
\toprule
\textbf{Scenario} & \textbf{Q-HAS Final L2} & \textbf{Classical Final L2} & \textbf{Winner} \\
\midrule
KH      & $2 \times 10^{-6}$ & $\sim 0$           & Classical \\
Tearing & $1.0 \times 10^{-5}$ & $3 \times 10^{-6}$ & Classical \\
Rotor   & $9.6 \times 10^{-4}$ & $5.2 \times 10^{-4}$ & Classical \\
OT      & $1.6 \times 10^{-4}$ & $\sim 0$           & Classical \\
\bottomrule
\end{tabular}
\end{table}

Both methods achieve excellent fidelity (all L2 errors $< 0.001$). The differences are practically negligible for KH and Tearing (L2 $< 10^{-5}$). Energy and enstrophy conservation track closely for both methods.

\subsection{Spatial Coherence}

\begin{figure}[H]
    \centering
    \includegraphics[width=\textwidth]{figures/fig3_spatial_coherence.png}
    \caption{Spatial Coherence analysis}
    \label{fig:spatial-coherence}
\end{figure}


\begin{table}[H]
\centering
\caption{Spatial coherence metrics (lower = more coherent). Classical AMR is consistently more spatially coherent.}
\label{tab:coherence}
\begin{tabular}{lcccc}
\toprule
 & \multicolumn{2}{c}{\textbf{Compactness (P/A)}} & \multicolumn{2}{c}{\textbf{Component Density}} \\
\textbf{Scenario} & Q-HAS & Classical & Q-HAS & Classical \\
\midrule
KH      & 0.048 & \textbf{0.040} & 0.055 & 0.055 \\
Tearing & 0.075 & \textbf{0.070} & 0.147 & \textbf{0.122} \\
Rotor   & 0.142 & \textbf{0.067} & 0.337 & \textbf{0.137} \\
OT      & 0.221 & \textbf{0.190} & 2.119 & \textbf{0.925} \\
\bottomrule
\end{tabular}
\end{table}

This is counterintuitive: the ZZ correlations were designed to enforce spatial consistency (ferromagnetic coupling), but in practice Q-HAS produces \textit{more} fragmented patch layouts. The perturbative corrections cause fragmentation because some cells get pushed above threshold while adjacent cells get pushed below. The ZZ coupling is too weak (suppressed by $\sigma = 0.023$) to impose coherence.

\subsection{Early Detection and Temporal Stability}

\begin{figure}[H]
    \centering
    \includegraphics[width=\textwidth]{figures/fig2_early_detection.png}
    \caption{Early Detection and Temporal Stability analysis}
    \label{fig:temporal}
\end{figure}


In Fig.(\ref{fig:temporal}), Q-HAS shows slightly better F1 than classical on KH and OT (Precision difference :$\Delta \approx 0.04$), driven by higher precision at the cost of marginally lower recall. However, classical is more temporally stable across all scenarios (IoU ranges $\Delta_{max} \approx 0.25$). The QAOA corrections introduce slight temporal instability in patch selections.

\subsection{Depth Analysis}

\begin{figure}[H]
    \centering
    \includegraphics[width=\textwidth]{figures/fig12_depth_analysis.png}
    \caption{Depth analysis}
    \label{fig:depth-analysis}
\end{figure}

In Fig.(\ref{fig:depth-analysis}), through depths 0--4, Q-HAS and Classical agree $>95\%$ of the time. Disagreements concentrate at the deepest level (depth 5), where agreement drops to $\sim 65\%$ on KH. This confirms the corrections are local and perturbative: the hierarchical structure is robust.

% ===========================
\chapter{Discussion of Results}
\label{chap:discussion}

\section{Discussion of Theoretical Results}

The gradient force (Equation~\eqref{eq:gradient-force}) connects the physical plasma state ($\theta$, $\phi$) to the variational optimisation direction.

\paragraph{Phase encoding and topological validation are coupled.}
The gradient force contains both a \textit{phase boost} factor ($\sin(\phi_q)$, growth rate) and a \textit{topological coherence} factor ($\prod_{m \neq q} \cos(\theta_m)$, spatial context). The algorithm simultaneously evaluates the dynamical tendency and spatial consistency of a potential instability.

\paragraph{Null-gradient proposition and symmetry breaking.}
The proof that $\nabla_{\Gamma} E = 0$ at $\vec{\Omega} = 0$ (Section~\ref{sec:null-gradient}) shows that structural parameters alone cannot initiate optimisation; the mixer must activate first. This justifies the linear-ramp strategy for $\vec{\Gamma}$.

\paragraph{Noise suppression through destructive interference.}
The coherence filter (Section~\ref{sec:coherence-filter}) distinguishes genuine instabilities from noise. In incoherent regions, the topological coherence factor averages to zero, suppressing the gradient and preventing false-positive refinement.

\section{Cross-Figure Synthesis}

The computational results collectively establish the following:

\begin{enumerate}
    \item \textbf{The QAOA circuit is functional.} Fig.~\ref{fig:qaoa-detail} proves non-zero corrections ($\delta$) exist at every BFS node. The corrections are structured (vary by scenario, depth, cell position) and are not random noise.

    \item \textbf{The corrections are perturbative, not transformative.} Fig.~\ref{fig:comprehensive} and Fig.~\ref{fig:statistical} show that practical outcomes (captured fraction, precision, recall) differ by at most $\pm 5\%$ between methods.

    \item \textbf{The only scenario with a genuine efficiency advantage is KH.} Q-HAS uses 10.9\% less compute with nearly the same capture quality, and this is statistically significant ($p = 0.008$). On all other scenarios, Q-HAS is either equivalent or worse.

    \item \textbf{ZZ correlations fragment rather than consolidate patches (Fig.~\ref{fig:spatial-coherence}).} Counter to design intent --- the ferromagnetic ZZ coupling was introduced to enforce spatial coherence --- the perturbative corrections produced at the trained $\sigma = 0.023$ are too weak to impose agreement and instead perturb decisions heterogeneously near the threshold, yielding more fragmented patches than the classical baseline. This places the coupling in a regime where it modifies decisions without being strong enough to coarsen them into coherent regions.

    \item \textbf{The root cause is scale.} At $2\times 2$ VQA resolution, only 4 cells are coupled, $\sigma = 0.023$ suppresses ZZ for most cells, $H(Z)$ dominates the Hamiltonian, and classical scoring is near-optimal at $Re = 400$. Disagreements concentrate at the deepest BFS level ($>95\%$ agreement through depths 0--4, dropping to $\sim 65\%$ at depth 5 on KH). With $\sigma = 0.023$, KH's ZZ coefficients are $10^4\times$ weaker than Z; only OT has non-trivial ZZ coupling. OT is the only scenario where Q-HAS wins on decision accuracy, confirming the topology-dependent advantage hypothesis.

    \item \textbf{Topology-rich regions show measurable Q-HAS advantage (Fig.~\ref{fig:topological}).} ZZ-driven flips are 76.5\% accurate for Tearing; ZZZZ-driven flips are 66.3\% accurate for OT. In topology-rich sub-regions, Q-HAS has $+3.8$pp advantage for Tearing and $+5.3$pp for OT. This is the strongest evidence that the Hamiltonian's spatial coupling terms detect real physics, even though the overall advantage is diluted because most of the domain is topologically smooth.
\end{enumerate}

\section{Identified Errors and Methodological Limitations}
\label{sec:identified-errors}

A careful analysis of Q-HAS's performance reveals several technical and methodological limitations that qualify the conclusions and delimit the scope of the claims made in this work.

\paragraph{Solver- and design-level issues.}
\begin{enumerate}
    \item \textbf{Failure of topological corrections on Harris tearing (Fig.~\ref{fig:topological}).} Q-HAS performs 97 \textit{incorrect} flips against only 50 \textit{correct} flips relative to ground truth on this scenario. The $ZZZZ$ X-point detection term is poorly calibrated at the current $2 \times 2$ VQA resolution: the plaquette coefficient couples four cells whose features are already aggregated by downsampling, compressing the dynamic range of $\det(\nabla \mathbf{B})$.
    \item \textbf{Dominance of the Z bias over spatial couplings (Fig.~\ref{fig:flip-analysis}).} At $w_{z,\text{frac}} = 10.40$ and the trained $\sigma = 0.023$, the activity-bias term dominates the cost Hamiltonian by roughly $10^3$--$10^4$ in coefficient magnitude on KH and Tearing. The majority of QAOA decisions are therefore near-copies of the classical score, and the $ZZ$/$ZZZZ$ terms act only perturbatively. This reduces the ``quantum'' component of the decision to a thin correction layer on top of the classical multi-indicator baseline.
    \item \textbf{Sensitivity to the ground-truth definition.} The hard-label rule (``refine'' $\equiv$ local L2 error above the $p = 75$ percentile of the domain-wide error distribution) is an arbitrary choice. The percentile sweep in Section~\ref{sec:robustness} (Phase 2B) shows that the LOSO collapse is stable across $p \in \{60, 70, 75, 80, 85, 90\}$, but within-scenario V1 results have not been re-run at alternative percentiles and may be more sensitive.
\end{enumerate}

\paragraph{Baseline and cost accounting.}
\begin{enumerate}
    \setcounter{enumi}{3}
    \item \textbf{The classical baseline is internal.} The ``classical AMR'' comparator is the 4-indicator threshold AMR (RMS of vorticity, divergence, $|J_z|$, L\"ohner) implemented in the same code base, \emph{not} a production MHD code such as FLASH, Athena or PLUTO. A reviewer should read the V1 results as ``Q-HAS vs.\ multi-indicator threshold AMR, this work'', not as ``Q-HAS vs.\ production AMR''. Comparison to a production code is left as further work.
    \item \textbf{Wall-clock cost is not measured.} The ``compute ratio'' metric reports the fraction of pixels refined, not CPU seconds. At $N = 256$ the QAOA calls on the Qiskit-Aer simulator dominate wall-clock time: a full Q-HAS trajectory is approximately two to three orders of magnitude slower than the classical baseline in wall time. The inequality~\eqref{eq:advantage-condition} of Section~\ref{sec:complexity} is therefore not satisfied in the present regime, even ignoring $S_{\text{quant}}/S_{\text{class}}$.
    \item \textbf{A pure machine-learning baseline is not included.} The random-split mean-field GBT achieves $F_1 = 0.986$ (Table~\ref{tab:ceiling-collapse}), which exceeds both Q-HAS and the classical multi-indicator baseline within-scenario. A head-to-head comparison with an ML-based AMR criterion is therefore warranted. The case for pursuing a Hamiltonian approach over such a classical ML baseline rests on (i) interpretability of the coefficient architecture, (ii) suitability for eventual execution on analogue Ising hardware, and (iii) structural results such as the ceiling proposition in Section~\ref{sec:ceiling-proposition}; we do not claim performance superiority over ML surrogates within-scenario.
\end{enumerate}

\paragraph{Quantum hardware and noise.}
\begin{enumerate}
    \setcounter{enumi}{6}
    \item \textbf{No hardware or noisy-simulator results are reported.} The implementation supports FakeFez (IBM noisy simulator) and real hardware via a command-line flag, but all results in Chapters~\ref{chap:results} and~\ref{chap:falsification} are from the Qiskit-Aer statevector simulator. The $0.66\%$ composite-loss advantage is therefore an idealised upper bound; on current NISQ hardware, two-qubit gate error rates of $10^{-2}$ together with readout error would likely consume the entire margin.
\end{enumerate}

\paragraph{Statistical and protocol issues.}
\begin{enumerate}
    \setcounter{enumi}{7}
    \item \textbf{The $n = 10$ seeds used for statistical testing is underpowered.} Effects reported with $|d| > 1$ (Tearing at $d = -3.41$ in particular) should be reproduced at $n \ge 30$ before being treated as reliable effect-size estimates. The $p$-values in Table~\ref{tab:statistical} are also not corrected for multiple comparisons across the four scenarios.
    \item \textbf{Only four scenarios enter the LOSO protocol.} With $N_{\text{folds}} = 4$, the mean LOSO $F_1 = 0.189 \pm 0.150$ is dominated by inter-scenario variability. Extending the fold set to include Lamb--Oseen, island coalescence, double tearing, magnetic twist and noisy uniform scenarios (already implemented in the code base but not used in LOSO) would stabilise the estimate.
\end{enumerate}

\section{Where the Potential Lies}

The quantum advantage is expected to grow along three axes. First, at higher VQA resolution ($4\times 4$ or $8\times 8$), there is more spatial structure for ZZ correlations to exploit; boundary correction potential scales as $\mathcal{O}(N)$ while total cells scale as $\mathcal{O}(N^2)$. Second, at higher Reynolds numbers ($Re/Rm \gg 800$), thinner current sheets, stronger turbulent cascades, and more reconnection events produce more borderline cells near the threshold and strengthen the dominant $\beta_{\text{curl}}$ signal (44\% fANOVA); classical scoring becomes noisier while ZZ correlations can ``fill in'' missing cells (retraining would be required). Third, scenarios with richer topological structure would amplify the localised advantage demonstrated in Fig.~\ref{fig:topological} into a global one.

% ==========================================
% CHAPTER 9: FALSIFICATION STUDY
% ==========================================
\chapter{Falsification Study: Cross-Scenario Ceiling Analysis}
\label{chap:falsification}

The V1 results (Chapter~\ref{chap:results}) established that Q-HAS delivers a marginal within-scenario advantage on some instability classes. This chapter asks a stronger question: can \textit{any} local Ising Hamiltonian in the 9-feature basis beat the classical multi-indicator baseline across \textit{unseen} MHD instability classes? The answer determines whether Q-HAS has an in-principle ceiling above classical AMR for scenario-universal deployment.

\section{Motivation and Research Question}
\label{sec:falsification-motivation}

The V1 evaluation (Chapter~\ref{chap:results}) was conducted per-scenario: the Hamiltonian was tuned on a given instability class and evaluated on the same class. This protocol measures \textit{fitting quality} but not \textit{generalisation}. In a production setting, the AMR steering mechanism must handle unseen instability types without retraining.

The falsification study reformulates the question structurally. Rather than testing one specific Hamiltonian $H$, we bound the F1 achievable by \textit{any} Hamiltonian in the family of local Ising models whose coefficients are measurable functions of the 9 per-site features used in this work: the classical multi-indicator score $s$, $|\mathbf{v}|^2$, $|\mathbf{B}|^2$, $|\omega_z|$, $|J_z|$, $|\nabla \mathbf{v}|^2$, $|\nabla \mathbf{B}|^2$, $\det(\nabla \mathbf{B})$, and $Re$. This family includes the V1 Hamiltonian, the parameter-free V2 Hamiltonian, and any future Ising-type cost function that a QAOA, simulated annealing, or D-Wave solver might optimise.

\section{V2 Hamiltonian and Experimental Design}
\label{sec:v2-hamiltonian}

\subsection{The Parameter-Free V2 Hamiltonian}

The V2 Hamiltonian eliminates all tunable coupling coefficients from V1. It retains only a single-qubit $Z$ bias proportional to the classical score, controlled by two parameters: $c_\text{bias}$ (bias strength) and $\text{thr}_\text{amr}$ (decision threshold). All multi-body terms ($ZZ$, $ZZZZ$) are dropped, isolating the question of whether the feature basis itself can support cross-scenario discrimination. This is deliberately minimal: if the mean-field (single-site) Hamiltonian cannot generalise, then adding couplings cannot help (as we prove formally in Section~\ref{sec:ceiling-proposition}).

\subsection{Ground-Truth Labelling}

Each cell is labelled \textit{hard} (refine) or \textit{easy} (coarsen) by comparing its local L2 error against a high-resolution DNS reference. A cell whose L2 error exceeds the $p$-th percentile (default $p = 75$) of the domain-wide error distribution is labelled hard. This produces a binary AMR mask $y_i \in \{0,1\}$ with positive rate $\approx 0.25$.

\subsection{Ceiling Estimation via Machine-Learning Surrogates}

Rather than testing one Hamiltonian at a time, we estimate the best F1 achievable by \textit{any} per-site decision function $f(\phi_i) \to \{0,1\}$ over the 9 features. Three classifiers of increasing capacity serve as surrogates for the supremum:
\begin{itemize}
    \item \textbf{Logistic regression (LR):} linear decision boundary,
    \item \textbf{Random forest (RF):} axis-aligned nonlinear boundary,
    \item \textbf{Histogram gradient-boosted trees (GBT):} smooth nonlinear boundary.
\end{itemize}
When all three converge to a similar F1, the empirical ceiling is close to the true supremum $F1^*_\text{MF}$. A neighbourhood (stencil) GBT using 45 features (the cell's own 9 features plus those of its 4 periodic neighbours) estimates the stencil ceiling $F1^*_\text{ST}$, which bounds Hamiltonians with nearest-neighbour couplings.

All probability thresholds are fitted on \textit{training} predictions to avoid data leakage.

\subsection{Evaluation Protocol: Leave-One-Scenario-Out (LOSO)}

The key methodological contribution is the LOSO cross-validation protocol. Four folds are defined, one per instability class (Orszag--Tang, Harris tearing, Kelvin--Helmholtz, MHD rotor). In each fold, the model trains on three scenarios and is evaluated on the held-out fourth. This tests whether the learned decision function generalises to unseen physics, not merely to unseen snapshots of known physics. Statistical significance is assessed via snapshot-level paired bootstrap ($B = 500$, percentile CIs).

All experiments use $N = 256$ grid resolution, $Re \in \{400, 800, 1200, 1600\}$, 4 scenarios, $\text{dim} = 4$ (coarsening factor), and seed $= 0$ unless otherwise noted.

% ------------------------------------------------------------------
\section{The Solver Is Not the Bottleneck}
\label{sec:solver-bottleneck}

Before testing the family of Hamiltonians, we establish that the solver quality is not the limiting factor.

\begin{table}[H]
\centering
\caption{F1 of the v2 Hamiltonian under different solvers, compared to the classical multi-indicator baseline. The gap between SA and classical proves that the Hamiltonian, not the solver, is the bottleneck.}
\label{tab:solver-comparison}
\begin{tabular}{lcc}
\toprule
\textbf{Method} & \textbf{F1} & \textbf{Source} \\
\midrule
Classical multi-indicator & 0.475 & Phase 11 \\
Simulated annealing on v2 $H$ & 0.336 & Phase 7 \\
QAOA (reps = 2, state-vector) & $\approx$ SA & Phase 5 \\
Exact diagonalisation (dim = 2) & $\approx$ SA & Phase 4 \\
\bottomrule
\end{tabular}
\end{table}

SA reaches the exact ground state of the v2 Hamiltonian and still loses to the classical indicator by 0.139 in F1. Since exact diagonalisation and QAOA converge to the same ground state, no deeper QAOA circuit or better optimiser can close this gap. The bottleneck is the \textit{Hamiltonian}, not the \textit{solver}.

% ------------------------------------------------------------------
\section{The V2 Hamiltonian Is Not Scenario-Universal}
\label{sec:cmaes-universality}

Phase~10 applies CMA-ES to optimise the two free parameters $(c_\text{bias}, \text{thr}_\text{amr})$ of the v2 Hamiltonian, both jointly across all scenarios and per-scenario independently.

\begin{itemize}
    \item \textbf{Joint training:} $\Delta$F1 vs classical $= -0.003$ (CMA-ES, $N = 256$).
    \item \textbf{Per-scenario optima:} $c_\text{bias}^* \in \{0.38,\ 2.50,\ 5.25,\ 99.43\}$, a $264\times$ spread. The MHD rotor optimum hits the upper bound of the search grid.
    \item \textbf{Per-scenario deltas:} $\Delta$F1 $\leq 0$ in all 4 scenarios (OT: $0.000$, Harris: $-0.021$, KH: $-0.005$, Rotor: $-0.071$).
\end{itemize}

The $264\times$ spread in $c_\text{bias}^*$ means no single parameter setting can serve all scenarios. The v2 Hamiltonian is not scenario-universal even with optimal tuning.

% ------------------------------------------------------------------
\section{Random-Split Ceiling vs LOSO Collapse}
\label{sec:ceiling-collapse}

The central result of the falsification study is the contrast between two evaluation protocols applied to the same data and classifiers.

\begin{table}[H]
\centering
\caption{Mean-field and stencil ceilings under random split vs LOSO. The random-split ceiling is scenario memorisation; under LOSO it collapses by $\approx 5\times$. All numbers at $N = 256$, 440 snapshots, 7040 cells.}
\label{tab:ceiling-collapse}
\begin{tabular}{lll}
\toprule
\textbf{Classifier} & \textbf{Split} & \textbf{F1 (95\% CI)} \\
\midrule
Classical multi-indicator & Random & 0.471 [0.423, 0.514] \\
Mean-field GBT (9 features) & Random & \textbf{0.986} [0.972, 0.995] \\
Stencil GBT (45 features) & Random & 0.990 [0.978, 0.998] \\
\midrule
Classical multi-indicator & LOSO ($n=4$) & 0.434 $\pm$ 0.148 \\
Mean-field GBT (9 features) & LOSO ($n=4$) & \textbf{0.189} $\pm$ 0.150 \\
Stencil GBT (45 features) & LOSO ($n=4$) & 0.215 $\pm$ 0.142 \\
Learned linear $H$ (phase 11c) & LOSO ($n=4$) & 0.391 \\
\bottomrule
\end{tabular}
\end{table}

The random-split GBT achieves $F1 = 0.986$, but this reflects \textit{scenario memorisation}: the model learns scenario-specific feature distributions from training snapshots and applies them to held-out snapshots of the same scenarios. Under LOSO, the same model collapses to $F1 = 0.189$, well below the classical baseline of 0.434.

The stencil GBT (45 features, covering each cell's nearest-neighbour context) achieves only $F1 = 0.215$ under LOSO, confirming that nearest-neighbour spatial couplings ($ZZ$, $ZZZZ$ terms) cannot rescue the collapse.

\subsection{Per-Fold Bootstrap Analysis}

Snapshot-level paired bootstrap ($B = 500$) provides per-fold confidence intervals on the difference $\Delta = F1_\text{site} - F1_\text{class}$:

\begin{table}[H]
\centering
\caption{Per-fold LOSO bootstrap results. 3/4 folds reject site $\geq$ classical at $p < 0.05$. Orszag--Tang is the only fold with a positive delta.}
\label{tab:loso-bootstrap}
\begin{tabular}{lccc}
\toprule
\textbf{Held-out} & $\bm{\Delta}$ \textbf{site--class} & \textbf{95\% CI} & $\bm{p}$\textbf{(site} $\geq$ \textbf{class)} \\
\midrule
Orszag--Tang & $+0.063$ & $[+0.043, +0.084]$ & 1.000 \\
Harris tearing & $-0.400$ & $[-0.463, -0.338]$ & 0.000 \\
Kelvin--Helmholtz & $-0.056$ & $[-0.067, -0.046]$ & 0.000 \\
MHD rotor & $-0.588$ & $[-0.633, -0.544]$ & 0.000 \\
\bottomrule
\end{tabular}
\end{table}

Harris and KH both reach $F1_\text{class} = 0.400$, which is the degenerate all-positive prediction floor: with a positive rate of 0.250, predicting all cells as positive yields $F1 = 2 \times 0.25 / (1 + 0.25) = 0.400$ (precision $= 0.25$, recall $= 1.0$). The classical threshold fails to discriminate on these held-out scenarios.

The stencil-vs-site gap has CI upper bound $= +0.009$ (phase 11H), confirming that the mean-field ceiling is tight: no Hamiltonian with nearest-neighbour couplings can exceed it by more than $\approx 0.01$.

% ------------------------------------------------------------------
\section{Robustness of the LOSO Collapse}
\label{sec:robustness}

Four independent tests confirm that the collapse is not an artefact of the labelling threshold, the classifier, the random seed, or missing scenario information.

\subsection{Percentile Sensitivity (Phase 2B)}

The L2-hard label threshold is swept across percentiles $p \in \{60, 70, 75, 80, 85, 90\}$. The LOSO delta ($F1_\text{site} - F1_\text{class}$) is negative at every percentile, ranging from $-0.163$ to $-0.306$. The collapse is robust to the definition of ``hard patch.''

\subsection{Scenario-Identity Ablation (Phase 11G)}

If the LOSO collapse were caused by the model lacking knowledge of \textit{which} scenario it faces, appending a one-hot scenario identifier should help.

\begin{table}[H]
\centering
\caption{Scenario-identity ablation under LOSO. Adding the correct or incorrect scenario identity has no effect ($\Delta \approx 0$).}
\label{tab:scenario-ablation}
\begin{tabular}{lc}
\toprule
\textbf{Feature set} & \textbf{LOSO mean F1} \\
\midrule
9 features only & 0.189 \\
9 features + correct one-hot ID & 0.188 \\
9 features + wrong one-hot ID & 0.196 \\
\bottomrule
\end{tabular}
\end{table}

Under LOSO, the held-out scenario's one-hot column is all-zero during training and all-one during validation, so the GBT cannot learn from a feature it never saw vary. The informative comparison is the ``wrong scenario'' row: the model \textit{does} see those one-hot columns vary in training, and shuffling scenario identity still yields $\Delta \approx 0$. Even the process-of-elimination signal does not help. The bottleneck is not ``which scenario am I in?'' but rather that the L2-hard label is not a per-snapshot function of the 9 features on a held-out instability class.

\subsection{Per-Scenario Specialisation (Phase 11D)}

Phase 11D trains a separate GBT on each scenario and evaluates it on all four, producing a $4 \times 4$ transfer matrix. Within-scenario (diagonal) F1 ranges from 0.90 to 0.98, confirming that the 9 features are sufficient \textit{within} a known instability class. Off-diagonal F1 ranges from 0.08 to 0.50, and the average off-diagonal-minus-classical delta is $-0.256$.

The binding constraint is MHD rotor: its within-scenario classical baseline is already $F1 = 0.925$, yet off-diagonal transfers into rotor average only 0.232. A scenario detector must achieve $\geq 97\%$ accuracy on rotor to break even with classical on that fold. Across all four folds, the worst-case break-even is 97\% (rotor) while the average break-even is 39\%. ``Train per scenario, switch at runtime'' is viable only with a near-perfect scenario detector for the hardest fold.

\subsection{Multi-Seed Validation (Phase 11F)}

Phase 11F reruns the full pipeline with 10 independent random seeds. Under LOSO, every fold gives identical F1 across all 10 seeds (to 3 significant figures), with per-fold $\sigma = 0.000$. The GBT fit is fully deterministic at this dataset size. The headline LOSO numbers ($0.189 \pm 0.150$) are not seed-limited; the $\pm 0.150$ is purely inter-scenario variability. Random-split across-seed $\sigma = 0.013$ confirms the phase~11H headline is reproducible within $\approx 2$ percentage points.

% ------------------------------------------------------------------
\section{V1 Pipeline Under LOSO}
\label{sec:v1-loso}

Phase~11E evaluates the V1 pipeline's input-side score (classical + $\psi$ temporal channel) under the V2 LOSO protocol with V1's best Optuna parameters (trial 85: $\beta = 0.5495$, $\text{thr}_\text{amr} = 0.3044$). Because V1's QAOA polishes the input by at most $\approx 0.66\%$ in combined decision score, this is a tight input-side proxy.

\begin{table}[H]
\centering
\caption{V1 input-side score under LOSO with snapshot-level paired bootstrap ($B = 500$, max-snaps/cfg $= 80$). 3/4 folds have CI strictly $> 0$.}
\label{tab:v1-loso}
\begin{tabular}{lccccc}
\toprule
\textbf{Held-out} & $\bm{F1_\text{v2,cls}}$ & $\bm{F1_\text{v1,cls}}$ & $\bm{F1_\text{v1+}\psi}$ & $\bm{\Delta}$ & \textbf{95\% CI} \\
\midrule
Orszag--Tang & 0.264 & 0.395 & 0.306 & $+0.042$ & $[+0.014, +0.070]$ \\
Harris tearing & 0.400 & 0.400 & 0.385 & $-0.015$ & $[-0.040, +0.006]$ \\
Kelvin--Helmholtz & 0.400 & 0.400 & 0.411 & $+0.011$ & $[+0.003, +0.023]$ \\
MHD rotor & 0.672 & 0.936 & 0.928 & $+0.256$ & $[+0.209, +0.302]$ \\
\midrule
\textbf{Mean} & 0.434 & 0.533 & 0.508 & $\bm{+0.074}$ & \\
\bottomrule
\end{tabular}
\end{table}

The $+0.074$ mean advantage decomposes as:
\begin{itemize}
    \item $F1_\text{v1,cls} - F1_\text{v2,cls} = +0.099$: the effect of \textbf{aggregation method}. Both V1 and V2 compute the identical 4-indicator classical score via \texttt{AngleMapper.classical\_score()} (vorticity, divergence, $|J_z|$, L\"ohner). The difference is downsampling: V1 averages within each coarse patch (\texttt{block\_avg}), while V2 takes the maximum (\texttt{block\_max}). Averaging preserves sub-patch variation that the threshold can exploit; the max operator saturates high-gradient patches and compresses the dynamic range.
    \item $F1_\text{v1+}\psi - F1_\text{v1,cls} = -0.025$: the $\psi$ temporal channel \textbf{hurts on average}. This is consistent with $\psi$ being tuned on within-scenario data (Optuna trial 85): it helps where it was trained but does not transfer cross-scenario.
\end{itemize}

V1's advantage under LOSO is entirely due to the aggregation method, not the temporal channel or QAOA polishing.

% ------------------------------------------------------------------
\section{Formal Ceiling Proposition}
\label{sec:ceiling-proposition}

The empirical results are elevated to a structural statement by the following proposition. We state first the mean-field version rigorously, then the stencil-range extension, and conclude with an operational corollary that makes explicit the assumption on the solver.

\paragraph{Setup.}
Fix a finite grid of $M$ cells $\{1, \dots, M\}$ with per-site feature vectors $\phi_i \in \mathbb{R}^9$ and binary ground-truth labels $y_i \in \{0, 1\}$. A \emph{per-site decision function} is any Borel-measurable map $f : \mathbb{R}^9 \to \{0,1\}$; we denote by $\mathcal{F}$ the class of such maps. For a decision vector $s \in \{0,1\}^M$, let $F_1(s, y)$ denote the standard $F_1$ score. Define the mean-field supremum
\[
    F_1^{\star}_{\mathrm{MF}} \;\equiv\; \sup_{f \in \mathcal{F}} F_1\!\bigl( (f(\phi_i))_{i=1}^M,\; y \bigr).
\]

\begin{tcolorbox}[title={Proposition 1 (Mean-field ceiling).}]
Let $H(s) = \sum_i h_i(\phi_i)\,s_i$ be any Hamiltonian on the cell Boolean lattice $\{0,1\}^M$ whose couplings vanish ($J_{ij} = K_{ijkl} = \cdots = 0$), where each $h_i : \mathbb{R}^9 \to \mathbb{R}$ is a Borel-measurable function of the per-site feature vector. Let $s^{\star}(H) \in \arg\min_{s \in \{0,1\}^M} H(s)$ be any exact ground state (ties broken arbitrarily). Then
\[
    F_1\!\bigl(s^{\star}(H),\; y\bigr) \;\le\; F_1^{\star}_{\mathrm{MF}}.
\]
\end{tcolorbox}

\noindent\emph{Proof.}
Because $H$ decomposes additively, the minimisation separates cell-wise:
$s^{\star}_i = \mathbf{1}\!\{h_i(\phi_i) < 0\}$ (with either $0$ or $1$ admissible when $h_i(\phi_i) = 0$). Define the per-site map $f_H(\phi) := \mathbf{1}\{h_i(\phi) < 0\}$; since $h_i$ is Borel-measurable, $f_H \in \mathcal{F}$. Then $s^{\star}_i = f_H(\phi_i)$ for all $i$, so $(s^{\star})$ is the image under a per-site decision function of the feature tensor. It is therefore one element of the class over which $F_1^{\star}_{\mathrm{MF}}$ is the supremum, which gives the claimed inequality. \hfill$\square$

\paragraph{Stencil extension.}
For a translation-invariant $k$-stencil, let $\phi_i^{(k)}$ denote the concatenation of $\phi_j$ for $j$ in a $(2k+1)^d$-neighbourhood of $i$. Define the stencil ceiling
\[
    F_1^{\star}_{\mathrm{ST}} \;\equiv\; \sup_{g : \mathbb{R}^{9(2k+1)^d} \to \{0,1\}} F_1\!\bigl( (g(\phi_i^{(k)}))_i,\; y \bigr).
\]

\begin{tcolorbox}[title={Proposition 2 (Stencil ceiling, conditional).}]
Let $H(s) = \sum_i h_i(\phi_i) s_i + \sum_{\langle i,j \rangle} J_{ij}(\phi_i, \phi_j)\,s_i s_j + \dots$ be a local Hamiltonian whose couplings have range at most $k$ (i.e.\ are functions of features within a $k$-neighbourhood). If its ground state $s^{\star}(H)$ is cell-wise determined by a function of the corresponding stencil features --- as happens, in particular, whenever the ground state is unique, and more generally under any tie-breaking rule that is a per-site function of the $k$-stencil --- then
\[
    F_1\!\bigl(s^{\star}(H),\; y\bigr) \;\le\; F_1^{\star}_{\mathrm{ST}}.
\]
\end{tcolorbox}

\noindent\emph{Remark.}
Unlike Proposition~1, Proposition~2 is not unconditional: with generic non-trivial couplings the ground state does not in general decompose as a per-site function of the stencil, because degeneracies can correlate cells at arbitrary distances. We therefore state Proposition~2 as \emph{conditional}, under the hypothesis that the ground state is a per-site function of the $k$-stencil. Empirically we have $F_1^{\star}_{\mathrm{ST}} - F_1^{\star}_{\mathrm{MF}} \le 0.009$ (stencil GBT vs.\ site GBT, Phase 11H), which tightly bounds the coupling-induced gain in this feature basis.

\paragraph{Operational corollary (solver assumption made explicit).}
Propositions~1 and 2 bound $F_1$ of the \emph{exact} ground state. Real solvers return approximate ground states:
\begin{itemize}
    \item QAOA at finite depth $p$ does \emph{not} in general converge to the ground state \cite{farhi2014,zhou2020qaoa}.
    \item Simulated annealing converges almost surely to the ground state only in the limit of infinitely slow logarithmic cooling \cite{geman1984sa}, which is not reached in finite runtime.
    \item Analogue (D-Wave-style) Ising solvers are subject to control noise and finite annealing times.
\end{itemize}
The propositions therefore do not literally exclude any of these algorithms from exceeding $F_1^{\star}_{\mathrm{MF}}$, since they may return states that are not ground states and happen to have higher $F_1$. However, such exceedance would be a property of the \emph{error of the solver}, not of the Hamiltonian, and it cannot be engineered reliably. The precise claim we make is:

\begin{quote}
\emph{For any local Ising Hamiltonian in the 9-feature family, if the solver is assumed to return the exact ground state (up to tie-breaking), then its $F_1$ is bounded by $F_1^{\star}_{\mathrm{ST}} \lesssim F_1^{\star}_{\mathrm{MF}} + 0.009$.}
\end{quote}

Combined with the empirical LOSO mean-field ceiling $F_1^{\star}_{\mathrm{MF}} = 0.189$ and the stencil gap $\le 0.026$, this gives an upper bound of $\approx 0.215$, well below the classical multi-indicator baseline $F_1 = 0.434$. We conclude that, \emph{under the solver-exactness assumption}, scenario-universal MHD adaptive refinement is not representable by a local Ising Hamiltonian on the 9 per-site features used here.

% ------------------------------------------------------------------
\section{V1 and V2 Reconciliation}
\label{sec:reconciliation}

V1 (Chapter~\ref{chap:results}) and V2 (this chapter) answer different questions:
\begin{itemize}
    \item \textbf{V1 asks:} does a well-tuned Hamiltonian with temporal encoding improve AMR decisions on scenarios it was trained on?
    \item \textbf{V2 asks:} can any local Ising Hamiltonian in the 9-feature family generalise to unseen instability classes?
\end{itemize}

\noindent Neither result implies the other. The per-scenario consistency is summarised in Table~\ref{tab:reconciliation}.

\begin{table}[H]
\centering
\caption{Scenario-by-scenario consistency between V1 (within-scenario) and V2 (LOSO).}
\label{tab:reconciliation}
\begin{tabular}{p{2.8cm}p{4.8cm}p{4.8cm}c}
\toprule
\textbf{Scenario} & \textbf{V1 result} & \textbf{V2 LOSO result} & \textbf{Consistent?} \\
\midrule
KH & $p = 0.008$, $-10.9\%$ compute & $F1_\text{site} = 0.344$ vs $F1_\text{class} = 0.400$ & Yes \\
OT & $+5.3$ pp decisions & $\Delta = +0.063$ (only positive fold) & Yes \\
Harris & 97 incorrect / 50 correct flips & $F1_\text{site} = 0.000$ & Yes \\
MHD rotor & No advantage reported & $\Delta = -0.588$ & Yes \\
\bottomrule
\end{tabular}
\end{table}

V1's positive KH $p$-value and V2's negative KH LOSO delta are not contradictions. V1 tested ``does the V1 Hamiltonian improve decision quality on KH when trained and tested on KH-like data?''. V2 tested ``does the 9-feature local-$H$ family generalise to KH when trained on the other three scenarios?''. The first measures fit quality on the training distribution; the second measures out-of-distribution generalisation.

\paragraph{What V1 covers that V2 does not.}
(i)~Temporal anticipation via $\psi$ (Fig.~2 protocol: 6 time points, early-detection recall/precision/IoU); (ii)~full Pareto over $\lambda_\text{cost}$; (iii)~coherence/compactness/component-density diagnostics; (iv)~per-scenario Optuna tuning. These V1 contributions are not invalidated by V2.

\paragraph{What V2 covers that V1 does not.}
(i)~Cross-scenario generalisation (LOSO); (ii)~model-family ceiling bounding every possible local Ising $H$; (iii)~sensitivity and significance testing on the collapse (percentile sweep, bootstrap, specialisation matrix); (iv)~DNS validation against published references.

\paragraph{The joint picture.}
V1 shows the Q-HAS Hamiltonian does the best it can under per-scenario tuning and delivers a real but small advantage on some scenarios. V2 shows the ceiling on any local Ising $H$ in the 9-feature family collapses under cross-scenario evaluation, so a scenario-universal $H$ in this family cannot exist regardless of how the Optuna search is extended. The binding constraint is \textbf{feature insufficiency on held-out MHD instability classes}: the 9 per-site physical features cannot reconstruct the L2-hard label on a class not in the training set, even when the class identity is given as input.

\section{Reproducibility and Data Availability}
\label{sec:reproducibility}

All experiments in this work are reproducible from the public repository \url{https://github.com/armandld/BA_Proj}. The repository is mirrored at a frozen Zenodo DOI (to be assigned upon arXiv submission). The following information is needed to reproduce the headline numbers:
\begin{itemize}
    \item \textbf{Software stack.} Python~3.11, Qiskit~1.x (exact pinned version in \texttt{requirements.txt}), NumPy~1.26, SciPy~1.11, Optuna~3.x. All quantum circuits are executed on the Qiskit-Aer statevector simulator; no hardware or noisy-simulator results are included in the headline tables.
    \item \textbf{Fixed physical parameters.} $\mathrm{Re} = \mathrm{Rm} = 400$, $\eta = 10^{-3}$, $c_s = 1.0$, $B_{z,\text{guide}} = 0.1$, grid $N = 256$ on the periodic domain $[0, 2\pi]^2$, classical time step adaptive under CFL with safety factor $0.4$.
    \item \textbf{Hyperparameters.} The five trainable quantum parameters at their Phase~1 best are $\beta = 9.94$, $\sigma = 0.023$, $\beta_{\text{curl}} = 4.27$, $\beta_{\text{xpoint}} = 2.39$, $w_{z,\text{frac}} = 10.40$. The four frozen parameters are $\text{threshold}_{\text{amr}} = 0.1496$ (from classical training), $\gamma_{\text{hydro}} = 2.0$, $\gamma_{\text{mag}} = 0.5$, $\kappa = 10.0$.
    \item \textbf{Random seeds.} All V1 figures use seed~$=0$ unless stated. The statistical validation (Table~\ref{tab:statistical}) aggregates seeds~$0$--$9$. V2 figures use seed~$=0$ throughout except for the multi-seed test of Phase~11F, which uses seeds~$0$--$9$.
    \item \textbf{Data.} Each Optuna trial uses a DNS reference trajectory precomputed once per scenario and cached; these trajectories are included in the Zenodo archive (approx.\ 1.2~GB per scenario).
\end{itemize}
Wall-clock runtimes reported in any future revision should be measured on a single machine and stated together with hardware specifications; the present work does not claim wall-clock advantage and does not report CPU seconds.

\chapter{Conclusions}
\label{chap:conclusions}

\section{Summary of V1 Results}

A consistent mapping from coarse-grained MHD observables to qubit states was constructed. The polar angle $\theta = 2 \arcsin(\sqrt{s})$ is derived from a classical multi-indicator score (RMS of vorticity, divergence, $|J_z|$, L\"ohner estimator), ensuring each qubit starts with the same probability as the classical AMR detector. The azimuthal phase $\phi$ encodes the stress flux growth rate. This dual encoding enables predictive detection: the QAOA refines the classical baseline using entanglement, detecting instabilities that are growing but below classical gradient thresholds.

A structured Hamiltonian was defined using the decoupled $f \times g \times T_{rcf}$ architecture, with individual terms corresponding to four physically distinct classes of MHD anomalies. The variational gradient analysis yielded two principal results: (i)~the null-gradient proposition, proving that optimisation requires mixer-induced symmetry breaking, and (ii)~explicit gradient expressions connecting physical plasma quantities to the direction and magnitude of qubit rotations.

The computational evaluation reveals a nuanced picture. Q-HAS achieves a $0.66\%$ advantage in composite training loss over $170+$ Optuna trials. After relabelling the one-sided permutation test to agree with the sign of Cohen's $d$ (see Table~\ref{tab:statistical}), only Kelvin--Helmholtz shows a genuine positive for Q-HAS ($p = 0.008$, $d = +0.90$); Harris Tearing shows a significant \emph{regression} ($p = 0.001$, $d = -3.41$ against Q-HAS); Rotor and Orszag--Tang are consistent with classical being better but do not reach significance at $n=10$ seeds. Classical AMR is more accurate overall on per-cell decisions (46.6\% vs 43.3\%), and Q-HAS produces more fragmented patches. The strongest evidence for the Hamiltonian's physics-detection mechanism comes from the topological attribution analysis: in topology-rich sub-regions, Q-HAS shows $+3.8$pp for Tearing and $+5.3$pp for OT, indicating that the spatial coupling terms detect real physics locally even when the global decision is net worse.

\section{Summary of V2 Results: The Falsification Study}

The V2 study establishes that the bottleneck limiting Q-HAS in the cross-scenario setting is not the quantum solver but the Hamiltonian family itself. Five claims are supported:

\begin{enumerate}
    \item \textbf{Solver is not the binding constraint.} Simulated annealing reaches the exact ground state of the V2 Hamiltonian (F1 = 0.336) and still loses to the classical multi-indicator baseline (F1 = 0.475). A QAOA circuit that returns the same ground state cannot do better; a QAOA that returns an approximate (non-ground) state could in principle do better by accident, but this cannot be engineered reliably (see Section~\ref{sec:ceiling-proposition}).
    \item \textbf{No scenario-universal parameters.} CMA-ES on $(c_\text{bias}, \text{thr}_\text{amr})$ yields a $264\times$ spread in optimal $c_\text{bias}$ across the four training scenarios, ruling out a single tuning that serves all of them.
    \item \textbf{Cross-scenario collapse.} The mean-field GBT ceiling drops from F1 $= 0.986$ (random split) to $0.189$ (LOSO), $5.2\times$ below the classical baseline ($0.434$). 3/4 LOSO folds reject ``site $\ge$ classical'' at $p < 0.05$.
    \item \textbf{Feature insufficiency.} Adding scenario identity as a one-hot input does not raise the LOSO ceiling ($\Delta \approx 0$), confirming that the 9 per-site features cannot reconstruct the L2-hard label on a held-out instability class even when class identity is known.
    \item \textbf{Temporal channel does not transfer.} V1's $\psi$ temporal channel degrades the LOSO signal by $-0.025$ on average; V1's within-scenario edge comes from the block-average aggregation method rather than from the quantum component.
\end{enumerate}

The formal ceiling proposition (Section~\ref{sec:ceiling-proposition}) gives, under the assumption that the solver returns the exact ground state, the bound $F_1(s^\star(H), y) \le F_1^\star_{\mathrm{ST}} \le F_1^\star_{\mathrm{MF}} + 0.009$ for every local Ising Hamiltonian whose coefficients are Borel-measurable functions of the 9-feature basis.

\section{Joint Conclusions}

V1 and V2 answer different questions and their results are consistent:
\begin{itemize}
    \item V1 shows: a well-tuned Hamiltonian with temporal encoding delivers a real but small within-scenario advantage on some scenarios, a temporal anticipation signal, and a known failure on Harris.
    \item V2 shows: the ceiling on any local Ising H in the 9-feature family collapses under cross-scenario evaluation, so a scenario-universal H cannot exist regardless of tuning.
\end{itemize}

The scope of the negative claim is:
\begin{quote}
\textit{No scenario-universal local Ising Hamiltonian in the 9 physical features tested here beats the classical multi-indicator baseline across held-out MHD instability classes.}
\end{quote}

The scope of the V1 positive claim, after V2, is:
\begin{quote}
\textit{Under per-scenario tuning and a temporal $\psi$ channel, the V1 Hamiltonian delivers a statistically significant within-scenario advantage on Kelvin--Helmholtz and an early-detection precision gain on KH and OT, while failing on Harris tearing and MHD rotor.}
\end{quote}

Neither claim falsifies the other.

\paragraph{Limitations.}
The principal scope restrictions of this work are the following. (i)~All quantum circuits are executed on a statevector simulator (Qiskit-Aer); no hardware or noisy-simulator results are reported in the headline tables, and the $0.66\%$ composite-loss advantage is therefore an idealised upper bound. (ii)~The classical comparator is a multi-indicator threshold AMR implemented in this code base, not a production MHD code; comparison to FLASH, Athena or PLUTO is outstanding. (iii)~The V2 ceiling analysis is restricted to 2D incompressible MHD, $\mathrm{Re} \le 1600$, 4 instability classes, and 9 per-site features. (iv)~Wall-clock runtime is not measured: ``compute ratio'' counts refined pixels, and the Q-HAS trajectory is in fact two to three orders of magnitude slower than the classical baseline in wall time on the simulator, which the asymptotic inequality~\eqref{eq:advantage-condition} of Section~\ref{sec:complexity} does not currently satisfy. (v)~Statistical testing uses $n = 10$ seeds without multiple-comparison correction; the $|d| = 3.41$ effect on Tearing should be reproduced at $n \ge 30$ before being taken at face value. (vi)~The Phase~11E evaluation is an input-side proxy for a full QAOA-over-V1 LOSO sweep, not that sweep itself.

\paragraph{Further work.}
Natural directions for raising the cross-scenario ceiling are: (i)~richer feature bases (Helmholtz decomposition of $\mathbf{v}$ and $\mathbf{B}$, Els\"asser variables, power-spectrum descriptors) that directly raise $F_1^\star_{\mathrm{MF}}$; (ii)~alternative coarse-graining and aggregation schemes (V1's block-avg already shifts the LOSO mean by $+0.099$ over block-max); (iii)~multi-step temporal protocols evaluated under LOSO rather than per-scenario; (iv)~scaling to $4 \times 4$ or $8 \times 8$ VQA grids, where $ZZ$ and $ZZZZ$ correlations have more spatial structure; (v)~extending the scenario set along several axes simultaneously (3D, higher $\mathrm{Re}$, reconnection with guide field, compressible regimes); (vi)~execution on real quantum hardware with explicit gate-noise budgets; (vii)~a direct head-to-head comparison with a classical ML-based AMR criterion (random forest or gradient-boosted tree) as a sanity baseline; (viii)~a diagrammatic reformulation of the $f \times g \times T_{\mathrm{rcf}}$ architecture that may connect the framework to the categorical-quantum programme \cite{nielsen2010,schuld2021effect} and make its compositional structure explicit.

Q-HAS demonstrates that structured quantum subroutines can serve as topology-aware diagnostics within classical simulation architectures. However, the falsification study reveals a fundamental limit: the 9 per-site physical features used in the current framework cannot reconstruct AMR decisions on unseen instability classes, regardless of how the Hamiltonian is constructed or the solver is chosen. Advancing beyond this limit requires enriching the feature basis before revisiting the quantum component.

\chapter*{Acknowledgements}
\addcontentsline{toc}{chapter}{Acknowledgements}

I would like to thank my supervisor, Dr.\ Brian Appelbe, for his guidance and constructive feedback throughout this project. To write this report and code the associated software, Claude AI (from Anthropic) was used in the following ways: (1) to help identify and resolve logic errors in the MHD solver modules by providing specific code snippets for debugging and optimization; and (2) to rephrase portions of the manuscript to improve English conciseness and clarity by providing my original drafts for stylistic editing. In accordance with departmental guidance, I have personally verified all AI-generated outputs for technical accuracy and physical correctness.

\begin{thebibliography}{99}

\bibitem{freidberg2014}
J.~P. Freidberg, \textit{Ideal MHD} (Cambridge University Press, 2014).

\bibitem{goedbloed2004}
J.~P.~H. Goedbloed and S.~Poedts, \textit{Principles of Magnetohydrodynamics} (Cambridge University Press, 2004).

\bibitem{wesson2011}
J.~Wesson, \textit{Tokamaks}, 4th ed. (Oxford University Press, 2011).

\bibitem{chandrasekhar1961}
S.~Chandrasekhar, \textit{Hydrodynamic and Hydromagnetic Stability} (Clarendon Press, Oxford, 1961).

\bibitem{furth1963}
H.~P. Furth, J.~Killeen, and M.~N. Rosenbluth, ``Finite-resistivity instabilities of a sheet pinch,'' \textit{Phys.\ Fluids} \textbf{6}, 459--484 (1963).

\bibitem{bateman1978}
G.~Bateman, \textit{MHD Instabilities} (MIT Press, 1978).

\bibitem{pope2000}
S.~B. Pope, \textit{Turbulent Flows} (Cambridge University Press, 2000).

\bibitem{berger1984}
M.~J. Berger and J.~Oliger, ``Adaptive mesh refinement for hyperbolic partial differential equations,'' \textit{J.\ Comput.\ Phys.} \textbf{53}, 484--512 (1984).

\bibitem{berger1989}
M.~J. Berger and P.~Colella, ``Local adaptive mesh refinement for shock hydrodynamics,'' \textit{J.\ Comput.\ Phys.} \textbf{82}, 64--84 (1989).

\bibitem{nielsen2010}
M.~A. Nielsen and I.~L. Chuang, \textit{Quantum Computation and Quantum Information}, 10th anniversary ed. (Cambridge University Press, 2010).

\bibitem{harrow2009}
A.~W. Harrow, A.~Hassidim, and S.~Lloyd, ``Quantum algorithm for linear systems of equations,'' \textit{Phys.\ Rev.\ Lett.} \textbf{103}, 150502 (2009).

\bibitem{feynman1982}
R.~P. Feynman, ``Simulating physics with computers,'' \textit{Int.\ J.\ Theor.\ Phys.} \textbf{21}, 467--488 (1982).

\bibitem{farhi2014}
E.~Farhi, J.~Goldstone, and S.~Gutmann, ``A quantum approximate optimization algorithm,'' arXiv:1411.4028 (2014).

\bibitem{preskill2018}
J.~Preskill, ``Quantum computing in the NISQ era and beyond,'' \textit{Quantum} \textbf{2}, 79 (2018).

\bibitem{lucas2014}
A.~Lucas, ``Ising formulations of many NP problems,'' \textit{Front.\ Phys.} \textbf{2}, 5 (2014).

\bibitem{mcclean2018}
J.~R. McClean, S.~Boixo, V.~N. Smelyanskiy, R.~Babbush, and H.~Neven, ``Barren plateaus in quantum neural network training landscapes,'' \textit{Nat.\ Commun.} \textbf{9}, 4812 (2018).

\bibitem{goodfellow2016}
I.~Goodfellow, Y.~Bengio, and A.~Courville, \textit{Deep Learning} (MIT Press, 2016).

\bibitem{powell1994}
M.~J.~D. Powell, ``A direct search optimization method that models the objective and constraint functions by linear interpolation,'' in \textit{Advances in Optimization and Numerical Analysis}, edited by S.~Gomez and J.-P.~Hennart (Springer, 1994), pp.~51--67.

\bibitem{akiba2019}
T.~Akiba, S.~Sano, T.~Yanase, T.~Ohta, and M.~Koyama, ``Optuna: A next-generation hyperparameter optimization framework,'' in \textit{Proc.\ 25th ACM SIGKDD} (2019), pp.~2623--2631.

\bibitem{qiskit2024}
Qiskit contributors, ``Qiskit: An open-source framework for quantum computing,'' \url{https://github.com/Qiskit/qiskit} (2024).

\bibitem{numpy2020}
C.~R. Harris \textit{et al.}, ``Array programming with NumPy,'' \textit{Nature} \textbf{585}, 357--362 (2020).

\bibitem{scipy2020}
P.~Virtanen \textit{et al.}, ``SciPy 1.0: fundamental algorithms for scientific computing in Python,'' \textit{Nat.\ Methods} \textbf{17}, 261--272 (2020).

\bibitem{chorin1968}
A.~J. Chorin, ``Numerical solution of the Navier--Stokes equations,'' \textit{Math.\ Comput.} \textbf{22}, 745--762 (1968).

\bibitem{orszag1979}
S.~A. Orszag and C.~M. Tang, ``Small-scale structure of two-dimensional magnetohydrodynamic turbulence,'' \textit{J.\ Fluid Mech.} \textbf{90}, 129--143 (1979).

%\bibitem{brandt1977}
%A.~Brandt, ``Multi-level adaptive solutions to boundary-value problems,'' \textit{Math.\ Comput.} \textbf{31}, 333--390 (1977).

\bibitem{fryxell2000}
B.~Fryxell \textit{et al.}, ``FLASH: An adaptive mesh hydrodynamics code for modeling astrophysical thermonuclear flashes,'' \textit{Astrophys.\ J.\ Suppl.\ Ser.} \textbf{131}, 273--334 (2000).

\bibitem{stone2008}
J.~M. Stone, T.~A. Gardiner, P.~Teuben, J.~F. Hawley, and J.~B. Simon, ``Athena: a new code for astrophysical MHD,'' \textit{Astrophys.\ J.\ Suppl.\ Ser.} \textbf{178}, 137--177 (2008).

\bibitem{solon2015}
A.~P. Solon \textit{et al.}, ``Pressure is not a state function for generic active fluids,'' \textit{Phys.\ Rev.\ E} \textbf{92}, 062111 (2015).

\bibitem{lohner1987}
R.~L\"ohner, ``An adaptive finite element scheme for transient problems in CFD,'' \textit{Comput.\ Methods Appl.\ Mech.\ Eng.} \textbf{61}, 323--338 (1987).

\bibitem{okubo1970}
A.~Okubo, ``Horizontal dispersion of floatable particles in the vicinity of velocity singularities such as convergences,'' \textit{Deep-Sea Res.} \textbf{17}, 445--454 (1970).

\bibitem{weiss1991}
J.~Weiss, ``The dynamics of enstrophy transfer in two-dimensional hydrodynamics,'' \textit{Physica D} \textbf{48}, 273--294 (1991).

\bibitem{mignone2007}
A.~Mignone, P.~Rossi, G.~Bodo, A.~Ferrari, and S.~Massaglia, ``PLUTO: A numerical code for computational astrophysics,'' \textit{Astrophys.\ J.\ Suppl.\ Ser.} \textbf{170}, 228--242 (2007).

% --- Related work additions (quantum algorithms for PDEs and plasma) ---

\bibitem{gaitan2020}
F.~Gaitan, ``Finding flows of a Navier--Stokes fluid through quantum computing,'' \textit{npj Quantum Information} \textbf{6}, 61 (2020).

\bibitem{liu2021carleman}
J.-P. Liu, H.~{\O}ie Kolden, H.~K. Krovi, N.~F. Loureiro, K.~Trivisa, and A.~M. Childs, ``Efficient quantum algorithm for dissipative nonlinear differential equations,'' \textit{Proc.\ Natl.\ Acad.\ Sci.\ USA} \textbf{118}, e2026805118 (2021).

\bibitem{childs2020pde}
A.~M. Childs, J.-P. Liu, and A.~Ostrander, ``High-precision quantum algorithms for partial differential equations,'' \textit{Quantum} \textbf{5}, 574 (2021).

\bibitem{dodin2021}
I.~Y. Dodin and E.~A. Startsev, ``On applications of quantum computing to plasma simulations,'' \textit{Phys.\ Plasmas} \textbf{28}, 092101 (2021).

\bibitem{martyn2021qsvt}
J.~M. Martyn, Z.~M. Rossi, A.~K. Tan, and I.~L. Chuang, ``Grand unification of quantum algorithms,'' \textit{PRX Quantum} \textbf{2}, 040203 (2021).

\bibitem{katesharbeck2019}
J.~Kates-Harbeck, A.~Svyatkovskiy, and W.~Tang, ``Predicting disruptive instabilities in controlled fusion plasmas through deep learning,'' \textit{Nature} \textbf{568}, 526--531 (2019).

% --- QAOA-related additions ---

\bibitem{hadfield2019}
S.~Hadfield, Z.~Wang, B.~O'Gorman, E.~G. Rieffel, D.~Venturelli, and R.~Biswas, ``From the quantum approximate optimization algorithm to a quantum alternating operator ansatz,'' \textit{Algorithms} \textbf{12}, 34 (2019).

\bibitem{egger2021warmstart}
D.~J. Egger, J.~Mare\v{c}ek, and S.~Woerner, ``Warm-starting quantum optimization,'' \textit{Quantum} \textbf{5}, 479 (2021).

\bibitem{cerezo2021cost}
M.~Cerezo, A.~Sone, T.~Volkoff, L.~Cincio, and P.~J. Coles, ``Cost function dependent barren plateaus in shallow parametrized quantum circuits,'' \textit{Nat.\ Commun.} \textbf{12}, 1791 (2021).

\bibitem{zhou2020qaoa}
L.~Zhou, S.-T. Wang, S.~Choi, H.~Pichler, and M.~D. Lukin, ``Quantum approximate optimization algorithm: performance, mechanism, and implementation on near-term devices,'' \textit{Phys.\ Rev.\ X} \textbf{10}, 021067 (2020).

\bibitem{wang2018qaoa}
Z.~Wang, S.~Hadfield, Z.~Jiang, and E.~G. Rieffel, ``Quantum approximate optimization algorithm for MaxCut: a fermionic view,'' \textit{Phys.\ Rev.\ A} \textbf{97}, 022304 (2018).

% --- Expressivity and learnability ---

\bibitem{schuld2021effect}
M.~Schuld, R.~Sweke, and J.~J. Meyer, ``Effect of data encoding on the expressive power of variational quantum-machine-learning models,'' \textit{Phys.\ Rev.\ A} \textbf{103}, 032430 (2021).

\bibitem{cerezo2022vqa}
M.~Cerezo, A.~Arrasmith, R.~Babbush \textit{et al.}, ``Variational quantum algorithms,'' \textit{Nat.\ Rev.\ Phys.} \textbf{3}, 625--644 (2021).

\bibitem{du2022learnability}
Y.~Du, M.-H. Hsieh, T.~Liu, S.~You, and D.~Tao, ``Learnability of quantum neural networks,'' \textit{PRX Quantum} \textbf{2}, 040337 (2021).

\bibitem{geman1984sa}
S.~Geman and D.~Geman, ``Stochastic relaxation, Gibbs distributions, and the Bayesian restoration of images,'' \textit{IEEE Trans.\ Pattern Anal.\ Mach.\ Intell.} \textbf{PAMI-6}, 721--741 (1984).

\bibitem{tothBxB2000}
G.~T\'oth, ``The $\nabla \!\cdot\! \mathbf{B} = 0$ constraint in shock-capturing magnetohydrodynamic codes,'' \textit{J.\ Comput.\ Phys.} \textbf{161}, 605--652 (2000).

\end{thebibliography}


\end{document}